% generated by video/tools/script_handout.py -- do not edit
\PPart{Part II: Foundations}{Part II: Foundations}
\PSlide{Slide 10}{Orbitals: the independent-electron picture}{Slide 10: Orbitals: the independent-electron picture}
\PLine{Solving the many-electron problem exactly is hopeless for all but the smallest systems, so the orbital model replaces the pair repulsion by an average field that each electron feels from all the others.}
\PLine{Each electron then obeys a one-electron equation: the Fock operator, f, acting on an orbital phi p, gives the orbital energy epsilon p times phi p.}
\PLine{In practice each molecular orbital is built as a linear combination of atom-centred basis functions, chi mu, with coefficients c mu p. This is the LCAO approximation: linear combination of atomic orbitals.}
\PLine{Filling the orbitals with electrons, at most two per orbital and with opposite spins, builds the configuration. An orbital that holds a single electron is called a \textbf{SOMO}, a singly occupied molecular orbital.}
\PSlide{Slide 11}{Box 1: “Much ado about orbitals”}{Slide 11: Box 1: “Much ado about orbitals”}
\PLine{The review has a box called Much ado about orbitals. It says: the reality is that an orbital is a mathematical construct, an approximate, non-unique solution to the Schrödinger equation.}
\PLine{Non-unique, because mixing the occupied orbitals among themselves by a unitary transformation leaves the determinant, and so every observable, unchanged.}
\PLine{Approximate, because orbitals are defined only within the independent-particle model.}
\PLine{In nineteen ninety nine a report in Nature claimed the direct observation of copper’s d orbitals, and started a debate: can an orbital be observed at all? What can be observed is the electron density, the probability of finding an electron at a point, and orbital occupations.}
\PLine{Keep this in mind. It returns later, when the order of \textbf{restricted open-shell} orbital energies turns out to depend on a convention.}
\PPart{Part III: Spin \& antisymmetry}{Part III: Spin \& antisymmetry}
\PSlide{Slide 12}{Swapping two electrons}{Slide 12: Swapping two electrons}
\PLine{Start with two electrons. Define the exchange operator, P one two, which swaps all the coordinates of electrons one and two: position and spin together, written x.}
\PLine{\textbf{Step one.} Swapping twice restores the original order, so P one two squared is the identity.}
\PLine{\textbf{Step two.} If Psi is an eigenfunction of P one two with eigenvalue lambda, applying the operator twice gives lambda squared times Psi, and by step one that must equal Psi.}
\PLine{\textbf{Step three.} \Take{So lambda squared is one, and lambda is plus one or minus one.}}
\PLine{\textbf{Step four.} The \textbf{Hamiltonian} treats all electrons alike, so it commutes with P one two. That means the energy eigenstates can be chosen with a definite exchange symmetry.}
\PLine{For electrons, nature picks minus one. The paper puts it this way: the wavefunction for a system of fermions must be \textbf{antisymmetric} with respect to the interchange of any two particles. Fermions are particles of half-integer spin, such as electrons.}
\PSlide{Slide 13}{Slater determinants}{Slide 13: Slater determinants}
\PLine{\textbf{Antisymmetry} is built in automatically by a \textbf{Slater determinant}. For two electrons in spin orbitals chi one and chi two, it is chi one of electron one times chi two of electron two, minus chi two of electron one times chi one of electron two.}
\PLine{The rows label electrons and the columns label spin orbitals, and the determinant mixes every assignment of electrons to orbitals with a sign.}
\PLine{Swapping the two electrons swaps two rows of the determinant, which flips its sign: the function is \textbf{antisymmetric}.}
\PLine{If both columns hold the same spin orbital, the determinant has two equal columns and vanishes. That is the \textbf{Pauli} exclusion principle.}
\PLine{With a factor of one over root two it is normalised: each product has norm one and the two products are orthogonal, so the determinant’s squared norm is two.}
\PLine{\Take{So Pauli’s rule of nineteen twenty five, that no two electrons have all four quantum numbers equal, is a consequence of antisymmetry inside the orbital model, not a separate law.}}
\PSlide{Slide 14}{Spin operators for two electrons}{Slide 14: Spin operators for two electrons}
\PLine{Now spin. For one electron of spin one half, s z gives plus one half on alpha, spin up, and minus one half on beta, spin down. The raising operator, s plus, turns beta into alpha and gives zero on alpha, and the lowering operator, s minus, turns alpha into beta and gives zero on beta.}
\PLine{For two electrons we add each electron’s operator: S z is s z of electron one plus s z of electron two, and the same for S plus and S minus.}
\PLine{A useful identity follows. Write S minus times S plus using S plus or minus equals S x plus or minus i S y.}
\PLine{Multiply out: S x squared plus S y squared, plus i times the commutator of S x and S y.}
\PLine{The commutator of S x and S y is i S z, and S x squared plus S y squared is S squared minus S z squared.}
\PLine{And i times i is minus one.}
\PLine{Rearranged: S squared equals S minus S plus, plus S z squared, plus S z. This form is handy because S plus gives zero on many simple spin functions.}
\PSlide{Slide 15}{$\hat S^2$ on the four spin products (1/2)}{Slide 15: S² on the four spin products (1/2)}
\PLine{Now apply S squared to the four products of one-electron spins. First, both electrons up: alpha one alpha two. S plus gives zero, because there is nothing left to raise.}
\PLine{S z gives one half plus one half, which is one.}
\PLine{\Take{So S squared gives zero from the first term, one from S z squared, and one from S z, two in total.}}
\PLine{The eigenvalue of S squared is S times S plus one, so two means S equals one, with M S equal to plus one.}
\PLine{Next, both down: beta one beta two. S plus can raise electron one or electron two, giving alpha one beta two plus beta one alpha two.}
\PLine{S minus then lowers whichever spin was raised, giving beta one beta two twice.}
\PLine{Here M S is minus one, so S z squared plus S z is one minus one, zero.}
\PLine{\Take{So S squared on beta one beta two is two times beta one beta two: again S equals one, now with M S equal to minus one.}}
\PSlide{Slide 16}{$\hat S^2$ on the four spin products (2/2)}{Slide 16: S² on the four spin products (2/2)}
\PLine{That leaves the two mixed products, alpha one beta two and beta one alpha two, both with M S equal to zero.}
\PLine{S plus on alpha one beta two can only raise electron two, giving alpha one alpha two. S minus then lowers either electron, giving beta one alpha two plus alpha one beta two.}
\PLine{S z gives one half minus one half, zero, so the S z terms contribute nothing.}
\PLine{\Take{So S squared on alpha one beta two gives alpha one beta two plus beta one alpha two.} It has turned the function into a different one, so alpha one beta two is not an eigenfunction.}
\PLine{By the same steps, S squared on beta one alpha two gives the same sum.}
\PLine{Now add and subtract the two results: S squared on the sum, or the difference, of the two products is one plus or minus one times that same combination.}
\PLine{The sum has eigenvalue two, so S is one: a \textbf{triplet}. The difference has eigenvalue zero, so S is zero: a \textbf{singlet}.}
\PSlide{Slide 17}{Triplet and singlet spin functions}{Slide 17: Triplet and singlet spin functions}
\PLine{Normalising the results gives the four two-electron spin states: three \textbf{triplet} functions with M S equal to plus one, zero and minus one, and one \textbf{singlet}.}
\PLine{The factor one over root two appears because the two products are orthonormal, so their sum has squared norm two.}
\PLine{All three triplet spin functions are symmetric when the electrons are swapped, while the singlet is \textbf{antisymmetric}.}
\PLine{The total wavefunction must be antisymmetric, so a triplet must have an antisymmetric spatial part, and a singlet a symmetric one. This pairing of space and spin decides everything that follows.}
\PLine{Test T one checked all of this exactly: S squared and S z were built as four by four matrices in \textbf{SymPy}, and every eigenvalue relation holds exactly.}
\PPart{Part IV: Two electrons, two orbitals}{Part IV: Two electrons, two orbitals}
\PSlide{Slide 18}{Six ways to place two electrons}{Slide 18: Six ways to place two electrons}
\PLine{Now the heart of the matter: two electrons in two orbitals, phi a below and phi b above. This is Figure one of the review, rebuilt. There are exactly six ways to place the electrons.}
\PLine{In arrangement one, both electrons sit in phi a with opposite spins. In arrangement two, both sit in phi b.}
\PLine{In three, phi a holds an up spin and phi b a down spin. Four is the same with the spins exchanged. All four of these have M S equal to zero.}
\PLine{In five, both spins point up, so M S is plus one, and in six both point down, so M S is minus one.}
\PLine{One and two are \textbf{closed shells}. Three and four are \textbf{open shells} with M S equal to zero. Five and six are \textbf{high spin}.}
\PLine{Counting confirms there are no others: there are four spin orbitals, phi a alpha, phi a beta, phi b alpha and phi b beta, and choosing two of four gives six determinants.}
\PSlide{Slide 19}{The paper’s Figure 1}{Slide 19: The paper’s Figure 1}
\PLine{Here is the review’s own Figure one, with its caption: various arrangements, or determinants, of two electrons in two molecular orbitals and the corresponding values of the secondary spin quantum number. That secondary spin quantum number is M S.}
\PLine{It matches our reconstruction exactly. Now let us write each determinant out in full.}
\PSlide{Slide 20}{Writing the determinants out}{Slide 20: Writing the determinants out}
\PLine{To keep the algebra readable, write A for phi a of electron one times phi b of electron two, and B for the same orbitals with the electrons swapped.}
\PLine{A determinant is chi one of electron one times chi two of electron two, minus the swapped product. For row one, the spin orbitals are phi a alpha and phi a beta, so both electrons share phi a and the spin part is alpha one beta two minus beta one alpha two. Row two is the same with phi b.}
\PLine{Row three uses phi a alpha and phi b beta. Written out, that is phi a of one alpha one, phi b of two beta two, minus phi b of one beta one, phi a of two alpha two: in shorthand, A alpha one beta two minus B beta one alpha two. Notice that in each product the spatial part and the spin part of the same electron travel together.}
\PLine{Row four, with the spins exchanged, is A beta one alpha two minus B alpha one beta two.}
\PLine{Rows five and six have both spins equal, so the spin part factors out: A minus B, times alpha one alpha two, or times beta one beta two.}
\PSlide{Slide 21}{$M_S=\pm1$: one determinant each}{Slide 21: MS=±1: one determinant each}
\PLine{Row five, normalised, is one over root two times the \textbf{antisymmetric} space function A minus B, times alpha one alpha two. That is a \textbf{triplet} spin function with M S equal to plus one. \Take{So a single determinant is already a product of a spatial part and a \textbf{spin eigenfunction}.}}
\PLine{Row six is the same spatial function with both spins down: the M S equals minus one component.}
\PLine{The review says exactly this: the triplet with M S equal to plus one is one over root two times phi a of one phi b of two minus phi b of one phi a of two, times alpha one alpha two, and these states can be represented as the single determinants five and six.}
\PLine{That is exactly right: the high-spin components of the triplet are single-reference.}
\PSlide{Slide 22}{$M_S=0$: iii and iv are not spin eigenfunctions}{Slide 22: MS=0: iii and iv are not spin eigenfunctions}
\PLine{Now the M S equals zero determinants. Row three is A times alpha one beta two, minus B times beta one alpha two.}
\PLine{But neither spin product is a \textbf{spin eigenfunction}. From the spin frames, S squared turns alpha one beta two into alpha one beta two plus beta one alpha two, which is not a multiple of alpha one beta two.}
\PLine{\Take{So determinant three is half \textbf{singlet} and half \textbf{triplet}, and only combinations of three and four can have a definite spin.}}
\PLine{Half and half is literal: alpha one beta two equals one over root two times the sum of theta one zero and theta zero zero, so the expectation value of S squared is one half times two plus one half times zero, which is one. Test T one b measures exactly one.}
\PLine{The question is: plus or minus? Which combination is the singlet?}
\PSlide{Slide 23}{The sum $D_{\text{iii} }+D_{\text{iv} }$, every term tracked}{Slide 23: The sum Diii+Div, every term tracked}
\PLine{Let us add the two determinants, keeping track of every term.}
\PLine{Expanding determinants three and four gives four products: A alpha one beta two, minus B beta one alpha two, plus A beta one alpha two, minus B alpha one beta two.}
\PLine{Collect the terms that multiply A, and the terms that multiply B.}
\PLine{Both brackets hold the same spin function, so it factors out: A minus B, times alpha one beta two plus beta one alpha two.}
\PLine{Now put A and B back. The spatial factor changes sign when the two electrons swap places, so it is \textbf{antisymmetric}. The spin factor is root two times theta one zero, the symmetric \textbf{triplet} function with M S equal to zero.}
\PLine{An antisymmetric space times a symmetric spin is a triplet. \Take{So the sum of determinants three and four is the M S equals zero component of the triplet, not the \textbf{open-shell singlet}.}}
\PSlide{Slide 24}{The difference $D_{\text{iv} }-D_{\text{iii} }$}{Slide 24: The difference Div-Diii}
\PLine{Now the difference: determinant four minus determinant three.}
\PLine{Four is A beta one alpha two minus B alpha one beta two; subtracting three gives minus A alpha one beta two plus B beta one alpha two.}
\PLine{Collect A and B: each multiplies beta one alpha two minus alpha one beta two.}
\PLine{\Take{So the common bracket factors out, now with A plus B.}}
\PLine{Put A and B back. The spatial part, A plus B, does not change when the electrons swap: it is symmetric. The spin part is minus root two times theta zero zero, the \textbf{antisymmetric} \textbf{singlet}. The overall minus sign of a state is irrelevant.}
\PLine{A symmetric space times an antisymmetric spin: the difference is the \textbf{open-shell singlet}.}
\PSlide{Slide 25}{Normalisation}{Slide 25: Normalisation}
\PLine{To normalise, compute the squared norm of A plus or minus B: the norm of A, plus the norm of B, plus or minus twice their overlap.}
\PLine{The norm of A splits into the product of the two orbital integrals, each one, because the orbitals are normalised.}
\PLine{The overlap of A and B splits into phi a with phi b, times phi b with phi a, and each of those is zero, because the orbitals are orthogonal.}
\PLine{\Take{So the squared norm is two, for both signs.}}
\PLine{The spin parts also have squared norm two, so each combination needs one over root two for space and one over root two for spin: a factor of one half overall.}
\PLine{The results: the \textbf{triplet} with M S equal to zero is one half times D three plus D four, and the \textbf{open-shell singlet} is one half times D four minus D three.}
\PSlide{Slide 26}{Now read the paper}{Slide 26: Now read the paper}
\PLine{Now read what the review prints on page four hundred forty one.}
\PLine{The first line gives the \textbf{open-shell singlet} in product form: one over root two times phi a of one phi b of two plus phi b of one phi a of two, times one over root two times beta one alpha two minus alpha one beta two. Then the text says it is a combination of two determinants, three and four. The determinants are written exactly as ours, and the prefactor, one over root two times one over root two, is our one half. Look closely at the sign between them.}
\PSlide{Slide 27}{Spotted it}{Slide 27: Spotted it}
\PLine{Here it is, marked.}
\PLine{In green, the product form: symmetric space times \textbf{antisymmetric} spin. That is the \textbf{open-shell singlet}, and it is correct.}
\PLine{In rose, the expansion: one half times determinant three plus determinant four. That plus sign makes it exactly the combination we just proved to be the \textbf{triplet}. The two printed lines contradict each other.}
\PLine{The expansion should read one half times determinant four minus determinant three.}
\PSlide{Slide 28}{Fluctuating multiplicity, live in print}{Slide 28: Fluctuating multiplicity, live in print}
\PLine{One character, a plus instead of a minus, turns a spin of zero into a spin of one.}
\PLine{The section of the review is titled Fluctuating ground-state \textbf{multiplicity}, and between two lines, the \textbf{open-shell singlet} fluctuates into a \textbf{triplet}.}
\PLine{It survived peer review and typesetting. Two referees with parallel spins, presumably: the \textbf{Fermi hole} kept them from ever meeting the error.}
\PSlide{Slide 29}{Checked three ways: SymPy, $\hat S^2$, Lean}{Slide 29: Checked three ways: SymPy, S², Lean}
\PLine{We checked this three independent ways. First, test T one expands the determinants symbolically in \textbf{SymPy}: the \textbf{triplet} equals one half of three plus four, the \textbf{singlet} equals one half of four minus three, and the paper’s version is not the singlet.}
\PLine{Second, test T one b shares no algebra with the first. It builds both wavefunctions as numerical tensors from random orthonormal orbitals, exactly as the paper prints the determinants, and applies the total spin operator directly. The paper’s combination gives an S squared of two, a triplet; the corrected one gives zero, a singlet; determinant three alone is half and half, with expectation one; and the paper’s own product form also gives zero.}
\PLine{Third, the \textbf{Lean} proof assistant checks it formally. The theorem sum is triplet proves that D three plus D four equals A minus B times the symmetric spin sum, by unfolding the definitions and ring arithmetic. A second theorem proves the printed plus version cannot be the singlet: it plugs in one counterexample and simplifies.}
\PSlide{Slide 30}{What the two-electron problem teaches}{Slide 30: What the two-electron problem teaches}
\PLine{To sum up the two-electron problem. There are three singlets: the two \textbf{closed shells}, one and two, and the \textbf{open-shell singlet}.}
\PLine{There is one \textbf{triplet}, with three M S components.}
\PLine{Its M S equals plus and minus one components are single determinants, five and six, while its M S equals zero component needs three plus four.}
\PLine{The open-shell singlet is irreducibly two determinants, whatever M S you choose: there is no single-determinant component to fall back on.}
\PLine{\Take{Hence the review’s conclusion: the triplet is accurately described by a single reference determinant, but open-shell singlets require a multiconfigurational, or multireference, treatment.}}
\PLine{That conclusion is right, despite the \Err{typo}.}
\PPart{Part V: Energies: J and K}{Part V: Energies: J and K}
\PSlide{Slide 31}{The integrals we need}{Slide 31: The integrals we need}
\PLine{We need two kinds of integral. The first is the one-electron integral h p q: orbital phi p, times the one-electron \textbf{Hamiltonian} acting on phi q, integrated over all space. When p equals q this is the energy of one electron alone in that orbital.}
\PLine{The second is the two-electron integral, written p q r s in round brackets. Electron one sits in the charge distribution phi p phi q, electron two in phi r phi s, and they repel as one over r one two. This is called chemists’ notation: the first pair always belongs to electron one.}
\PLine{Two special cases matter. The \textbf{Coulomb integral} J p q is p p q q: two ordinary charge clouds, phi p squared and phi q squared, repelling each other. That is classical electrostatics.}
\PLine{The \textbf{exchange integral K} p q is p q p q. Here the charge distribution is the overlap density, phi p times phi q, repelling itself. It has no classical analogue. It appears only because the wavefunction is \textbf{antisymmetric}, through the swapped product B, and it will carry the whole \textbf{singlet} \textbf{triplet} splitting.}
\PSlide{Slide 32}{Two elementary integrals}{Slide 32: Two elementary integrals}
\PLine{Every energy in this chapter is built from three elementary matrix elements between orbital products. Assume real orthonormal orbitals.}
\PLine{First, the one-electron operator for electron one. It acts only on electron one’s coordinates, so the integral over electron one gives h p r, and the integral over electron two is just the overlap of phi q with phi s, which is delta q s.}
\PLine{Second, the operator for electron two: the same with the roles swapped, delta p r times h q s.}
\PLine{Third, the repulsion. One over r one two depends on both electrons at once, so nothing separates, and the result is the two-electron integral p r q s.}
\PLine{The bookkeeping rule: electron one carries the indices p and r, from the bra and the ket; electron two carries q and s.}
\PSlide{Slide 33}{Triplet energy, step by step (1/2)}{Slide 33: Triplet energy, step by step (1/2)}
\PLine{Now compute the \textbf{triplet} energy. The triplet is one over root two times A minus B, times a triplet spin function. The \textbf{Hamiltonian} contains no spin, so the energy is the spatial expectation value times the spin overlap, which is one.}
\PLine{Start with the one-electron part. Squaring one over root two gives a factor of one half, and there are four terms: A with A, B with B, and minus the two cross terms.}
\PLine{A with A: electron one carries a and a, electron two b and b, so the result is h a a times one, plus one times h b b: h a plus h b.}
\PLine{B with B is the same with the orbitals swapped.}
\PLine{The cross term A with B: each piece contains an overlap of a with b, which is zero, so orthogonality kills it.}
\PLine{\Take{So the one-electron part is one half of h a plus h b, twice: just h a plus h b.}}
\PSlide{Slide 34}{Triplet energy, step by step (2/2)}{Slide 34: Triplet energy, step by step (2/2)}
\PLine{Now the two-electron part: the same four terms, with one over r one two.}
\PLine{A with A: electron one carries a and a, electron two b and b, giving a a b b: the \textbf{Coulomb integral} J a b, two charge clouds repelling.}
\PLine{B with B gives the same clouds with the labels swapped: again J a b.}
\PLine{The cross term A with B: electron one carries a and b, electron two b and a, giving a b b a, which for real orbitals equals the \textbf{exchange integral K} a b. Each electron now sits in an overlap density.}
\PLine{Adding up with the factor one half and the minus sign from A minus B gives J a b minus K a b.}
\PLine{\Take{So the \textbf{triplet} energy is h a plus h b plus J a b minus K a b.} The cross terms carry the sign of the combination, and that sign is the whole story.}
\PSlide{Slide 35}{Open-shell singlet and the gap}{Slide 35: Open-shell singlet and the gap}
\PLine{The \textbf{open-shell singlet} is one over root two times A plus B, times the \textbf{singlet} spin function. The steps are identical, except that the cross terms now enter with a plus sign.}
\PLine{\Take{So the two-electron part is one half of J plus J plus K plus K, and the singlet energy is h a plus h b plus J a b plus K a b.}}
\PLine{Subtract the \textbf{triplet} energy: h and J cancel, and the singlet lies above the triplet by exactly two K a b.}
\PLine{With the same orbitals in both states, the whole singlet triplet splitting is carried by the \textbf{exchange integral}.}
\PSlide{Slide 36}{The closed shells and their couplings}{Slide 36: The closed shells and their couplings}
\PLine{Now the \textbf{closed shells}. With both electrons in phi a, the energy is two h a, one for each electron, plus a single repulsion, J a a, of the cloud phi a squared with itself. The same for phi b.}
\PLine{Between the two closed shells, the one-electron terms vanish by orthogonality, and the repulsion gives phi a phi a with phi b phi b, which is a b a b: the \textbf{exchange integral K} a b again.}
\PLine{The coupling between \textbf{closed shell} one and the \textbf{open-shell singlet} is one over root two times the sum of the matrix elements with A and with B.}
\PLine{For the one-electron part with A: the h of one term gives h a a times the overlap of a with b, zero; the h of two term gives the overlap of a with a, one, times h a b.}
\PLine{The repulsion with A gives a a a b, and B gives the same.}
\PLine{Altogether: root two times h a b plus a a a b. States of different total spin never couple, because their spin parts are orthogonal, so the \textbf{triplet} stays on its own.}
\PSlide{Slide 37}{The full two-electron, two-orbital Hamiltonian}{Slide 37: The full two-electron, two-orbital Hamiltonian}
\PLine{Collecting everything: in the basis of the three singlets, the two \textbf{closed shells} and the \textbf{open-shell singlet}, the \textbf{Hamiltonian} is a three by three matrix, and the \textbf{triplet} sits on its own.}
\PLine{The diagonal holds the three energies we derived. The closed shells couple through K a b, and each couples to the open-shell singlet through root two times h a b plus a mixed integral. Diagonalising this matrix gives the three exact \textbf{singlet} energies of the model.}
\PLine{If the two orbitals have different symmetry, like the sigma g and sigma u orbitals of H two, the mixed integrals vanish and the open-shell singlet decouples from the closed shells.}
\PLine{Test T two checked this by brute force: five random sets of integrals, built from a positive-definite kernel, were diagonalised with \textbf{PySCF}’s \textbf{full configuration interaction}, and the singlet and triplet roots match these closed forms to two times ten to the minus fourteen hartree. That is the limit of double-precision arithmetic.}
\PSlide{Slide 38}{Hund’s rule at frozen orbitals}{Slide 38: Hund’s rule at frozen orbitals}
\PLine{The review states \textbf{Hund’s rule} in these words: of the different multiplets resulting from different configurations of electrons in degenerate orbitals of an atom, those with greatest \textbf{multiplicity} have the lowest energy.}
\PLine{With the same orbitals in both states, our result reproduces it: the \textbf{triplet} lies below the \textbf{singlet} by two K, provided K is positive.}
\PLine{Two things remain to be proved. First, that K is never negative, which comes next. Second, whether using the same orbitals in both states is legitimate. We will see that it is not, and that the true reason for Hund’s rule is different.}
\PPart{Part VI: Fermi holes}{Part VI: Fermi holes}
\PSlide{Slide 39}{Where can the two electrons be?}{Slide 39: Where can the two electrons be?}
\PLine{Where can the two electrons be? The pair density, P of r one and r two, is the probability density of finding one electron at r one and another at r two. It is N times N minus one, times the wavefunction squared, summed over the spins and integrated over all the other electrons.}
\PLine{For our two electrons, after summing over spin, the \textbf{triplet}’s pair density is the square of the \textbf{antisymmetric} spatial function, A minus B. The \textbf{open-shell singlet}’s is the square of A plus B.}
\PLine{The prefactors cancel: N times N minus one is two, which cancels the one half from normalisation, and the spin functions are normalised, so their sum over spins gives one.}
\PLine{Now put both electrons at the same point. A and B become equal, so the triplet’s pair density is exactly zero, while the \textbf{singlet}’s is four phi a squared phi b squared.}
\PLine{This is the \textbf{Fermi hole}: in the triplet, the two electrons are never found at the same point, for every M S. It is exact, with no approximation beyond the form of the wavefunction.}
\PLine{The Coulomb hole, caused by the repulsion itself, is absent at this level. It needs correlation beyond a single configuration, which the review calls \textbf{dynamic correlation}.}
\PSlide{Slide 40}{Seeing the Fermi hole}{Slide 40: Seeing the Fermi hole}
\PLine{To see the \textbf{Fermi hole}, put two electrons in a one-dimensional box, in the two lowest orbitals, root two sine pi x and root two sine two pi x. Here is the \textbf{triplet}’s pair density: the brightness at each point is the probability of finding electron one at x one and electron two at x two.}
\PLine{Along the dashed diagonal, where x one equals x two, it is dark all the way: the two electrons are never at the same place.}
\PLine{The \textbf{open-shell singlet} is different. Its pair density is bright along the diagonal: the two electrons are often found together.}
\PLine{\Take{So for the same two orbitals, the triplet keeps its electrons apart and the \textbf{singlet} does not.} \Take{That is why the triplet’s electron repulsion carries a minus K.}}
\PSlide{Slide 41}{$K_{ab}\ge0$: a three-line proof}{Slide 41: Kab≥0: a three-line proof}
\PLine{Now prove that K is never negative. Define the overlap density rho as phi a times phi b, which is real. Then K is the self repulsion of rho: the double integral of rho at one point times rho at another, over their distance.}
\PLine{\textbf{Step one.} Write one over r as a sum of plane waves: an integral over all wave vectors k of four pi over k squared times e to the i k dot r. Every plane wave enters with a positive weight.}
\PLine{\textbf{Step two.} Insert this with r equal to r one minus r two. The exponential splits into a factor for electron one and a factor for electron two, so the two spatial integrals separate. They are the Fourier transform of rho and its complex conjugate.}
\PLine{\textbf{Step three.} Because rho is real, the two factors are complex conjugates, and their product is the squared modulus of the Fourier transform.}
\PLine{\Take{So K is a positive weight times a square, integrated: it can never be negative.}}
\PLine{It is zero only if phi a times phi b vanishes everywhere, meaning the two orbitals share no region of space. That idea returns later as disjointness.}
\PSlide{Slide 42}{$J_{ab}\ge K_{ab}$}{Slide 42: Jab≥Kab}
\PLine{One more inequality. Take f to be the \textbf{antisymmetric} function A minus B.}
\PLine{Its square is never negative and one over r one two is positive, so the integral of f squared over r one two is at least zero.}
\PLine{Expand f squared as A squared plus B squared minus two A B. The A squared and B squared integrals are each J, and the A B integral is K, so the result is J plus J minus K minus K, which is two times J minus K.}
\PLine{\Take{So J a b is at least K a b, which is at least zero.}}
\PLine{The \textbf{triplet}’s repulsion, J minus K, is never negative, and it lies below the \textbf{singlet}’s J plus K.}
\PLine{Test T three checked this on real molecular orbitals: for every pair of canonical Hartree Fock orbitals of water in the six thirty one G basis, the smallest K is one point one times ten to the minus three hartree, positive, and J is at least K for every pair.}
\PLine{\Take{So with \textbf{frozen orbitals}, the triplet energy really is the lower of the two.}}
\PPart{Part VII: Hund’s rule}{Part VII: Hund’s rule}
\PSlide{Slide 43}{Two explanations of the same rule}{Slide 43: Two explanations of the same rule}
\PLine{\textbf{Hund’s rule} says the \textbf{triplet} is lower, but why? The review gives two explanations. Slater’s original explanation was less electron electron repulsion in the triplet state. However, the review continues, the electron electron repulsion is actually greater in the high-spin state, and an alternative explanation, the \textbf{less screening} rule, has been offered.}
\PLine{Read with \textbf{frozen orbitals}, Slater’s argument is exactly our result: the repulsion changes by minus two K, while the kinetic energy and the nuclear attraction do not change. The triplet wins by lower repulsion.}
\PLine{In the less screening picture the orbitals relax. The repulsion goes up, the nuclear attraction becomes much more negative, and the kinetic energy rises. The triplet wins by stronger nuclear attraction.}
\PLine{Which is right? The \textbf{virial theorem}, together with the computed sign of the change in repulsion, decides.}
\PSlide{Slide 44}{The virial theorem by scaling (1/2)}{Slide 44: The virial theorem by scaling (1/2)}
\PLine{The \textbf{virial theorem} decides between them, and it can be derived by scaling. Take an atom with the nucleus at the origin, and a wavefunction that is exact, or at least optimised with respect to scaling.}
\PLine{\textbf{Step one.} Define a scaled wavefunction, Psi eta: multiply every electron coordinate by eta, and put a prefactor eta to the three N over two in front. Eta greater than one squeezes the whole electron cloud, eta less than one stretches it.}
\PLine{\textbf{Step two.} Its norm: the prefactor appears squared, giving eta to the three N times the integral of Psi at eta r, squared.}
\PLine{Substitute s equals eta r. The volume element is eta to the minus three N times d s, which cancels the prefactor, and the norm is one.}
\PLine{\Take{So the scaled function is still normalised and is a legal trial function for every eta.} The prefactor exists only to keep the norm one: three coordinates per electron, N electrons.}
\PSlide{Slide 45}{The virial theorem by scaling (2/2)}{Slide 45: The virial theorem by scaling (2/2)}
\PLine{\textbf{Step three.} The kinetic energy contains two derivatives, and each derivative with respect to r is eta times a derivative with respect to s. \Take{So the kinetic energy of the scaled function is eta squared times the original.}}
\PLine{\textbf{Step four.} Every potential term is one over a distance, and one over r is eta times one over s. \Take{So the potential energy scales with one power of eta.}}
\PLine{\textbf{Step five.} Add them: the energy as a function of eta is eta squared times T plus eta times V.}
\PLine{\textbf{Step six.} An exact state is stationary under any variation, including scaling, so the derivative at eta equal to one vanishes: two T plus V equals zero.}
\PLine{That is the \textbf{virial theorem}. The total energy equals minus the kinetic energy, and also one half of the potential energy. It is proved in \textbf{Lean} as the theorem scaling derivative, and for molecules the same holds at equilibrium geometries.}
\PSlide{Slide 46}{Frozen orbitals cannot be the mechanism}{Slide 46: Frozen orbitals cannot be the mechanism}
\PLine{Apply the \textbf{virial theorem} to both exact states and subtract. The energy difference equals minus the kinetic energy difference, and also one half of the change in nuclear attraction plus the change in repulsion.}
\PLine{Now take the frozen-orbital model. If both states use the same orbitals, the kinetic energy is the same, because it is a one-electron operator summed over the same occupied orbitals.}
\PLine{Then the virial theorem forces the energy difference to be zero. But the same model says the \textbf{triplet} is lower by two K. That is a contradiction.}
\PLine{\Take{So \textbf{frozen orbitals} violate the virial theorem.} They cannot tell us which energy component actually pays for the splitting.}
\PLine{Hongo and coworkers make the same point: this assumption leaves the kinetic energy unchanged, which cannot be correct, as each eigenstate must satisfy the virial theorem.}
\PSlide{Slide 47}{Less screening, as inequalities}{Slide 47: Less screening, as inequalities}
\PLine{Now turn \textbf{less screening} into inequalities. Assume three things: the \textbf{triplet} is lower, both states obey the \textbf{virial theorem}, and the repulsion is larger in the triplet, which is what accurate calculations show.}
\PLine{\textbf{Step one.} From E equals minus T, the kinetic energy change is minus the energy change, which is positive: the triplet’s electrons are more compressed.}
\PLine{\textbf{Step two.} From E equals one half V, the potential energy change is twice the energy change, which is negative.}
\PLine{\textbf{Step three.} The potential energy is attraction plus repulsion, so the change in attraction is two delta E minus the change in repulsion. Since the repulsion went up, this is even more negative than two delta E.}
\PLine{\textbf{Step four.} \Take{So the gain in nuclear attraction is larger than the loss from extra repulsion.}}
\PLine{The nucleus wins: stronger attraction over-compensates stronger repulsion. This chain of inequalities is proved in \textbf{Lean} as the theorem less screening.}
\PSlide{Slide 48}{The physical picture}{Slide 48: The physical picture}
\PLine{Here is the physical picture. In the \textbf{singlet}, the two electrons may come close to each other. Each one spends part of its time between the nucleus and the other electron, screening the nuclear charge, so the cloud stays large.}
\PLine{In the \textbf{triplet}, the \textbf{Fermi hole} keeps the two electrons apart. Each screens the nucleus less, so both feel more of its charge and the cloud contracts towards the nucleus.}
\PLine{The contraction lowers the nuclear attraction energy, raises the kinetic energy, and also raises the repulsion, because the electrons are closer on average. All three signs match the virial analysis.}
\PSlide{Slide 49}{Test 1: helium 1s2s, near-exact}{Slide 49: Test 1: helium 1s2s, near-exact}
\PLine{Test T four computes this exactly for helium’s one s two s states: the \textbf{triplet}, two three S, and the \textbf{singlet}, two one S, in a large Gaussian basis: the triplet within two hundredths of a millihartree of the exact value, and the singlet within about a tenth of a millihartree.}
\PLine{The triplet is lower by zero point zero two nine three hartree. Both states satisfy the \textbf{virial theorem} to six decimal places.}
\PLine{The kinetic energy goes up by exactly the same amount, as the virial theorem requires.}
\PLine{The nuclear attraction drops by zero point zero seven seven hartree, while the repulsion goes up, by zero point zero one eight five.}
\PLine{\Take{So Slater’s mechanism, lower repulsion, fails: the repulsion is higher in the triplet, and the stronger nuclear attraction pays for everything.} As a check, one half of the two potential changes adds up to minus zero point zero two nine three, the energy difference.}
\PSlide{Slide 50}{Test 2: the carbon atom, from the literature}{Slide 50: Test 2: the carbon atom, from the literature}
\PLine{Test two is the carbon atom, from the literature. Hongo and coworkers computed the lowest \textbf{triplet} and \textbf{singlet} terms of carbon, three P and one D, with Hartree Fock and with diffusion quantum Monte Carlo, which captures nearly all of the correlation. We read their Table one from the full text and computed the differences.}
\PLine{The pattern is the same as in helium: the repulsion is higher in the triplet, the attraction is much stronger, and the kinetic energy is higher.}
\PLine{\Take{So the review’s account, \textbf{less screening}, with the repulsion greater in the high-spin state, agrees with our helium result and with carbon.} Verified.}
\PSlide{Slide 51}{A side claim: same-spin correlation}{Slide 51: A side claim: same-spin correlation}
\PLine{A side claim in the review: the correlation energy for a pair of same-spin electrons is roughly an order of magnitude smaller than for an opposite-spin pair. The correlation energy is the exact energy minus the Hartree Fock energy.}
\PLine{Test T fifteen measured this ratio per pair. The dashed line marks one tenth, an order of magnitude.}
\PLine{For helium, comparing the \textbf{triplet}, two three S, with one same-spin pair, to the ground state, with one opposite-spin pair, the ratio is zero point zero two three. The claim holds.}
\PLine{But for the valence shells of ethylene, nitrogen and water, second-order perturbation theory gives zero point three to zero point four per pair. Compare the spin-component-scaled M P two method, which scales the two parts by one third and six fifths.}
\PLine{\Take{So the claim is true for compact two-electron systems, and overstated in general.}}
\PPart{Part VIII: Static correlation}{Part VIII: Static correlation}
\PSlide{Slide 52}{Static versus dynamic correlation}{Slide 52: Static versus dynamic correlation}
\PLine{The next mechanism is \textbf{static correlation}. The review says: the weaker the interaction between the two nonbonding molecular orbitals, the stronger the static correlation, and points to its Figure two.}
\PLine{That figure shows three cases where the orbital interaction weakens: breaking the bond of H two, twisting ethylene, and stretching a metal ligand bond.}
\PLine{Two kinds of correlation are distinguished. \textbf{Dynamic correlation} is electrons dodging each other within one configuration. Static correlation means that several configurations have comparable weight, so no single determinant is a good starting point.}
\PSlide{Slide 53}{H$_2$ in a minimal basis}{Slide 53: H2 in a minimal basis}
\PLine{Take the simplest model of H two: one hydrogen one s function on each atom, chi A and chi B, with overlap S.}
\PLine{The bonding orbital, sigma g, is their normalised sum, and the antibonding orbital, sigma u, their normalised difference.}
\PLine{The normalisation factors follow from the norm of chi A plus or minus chi B: one plus one plus or minus two S.}
\PLine{Under inversion through the bond centre, sigma g is even and sigma u is odd. Any integral containing an odd number of sigma u functions has an odd integrand and integrates to zero. \Take{So the mixed one-electron integral, and the integrals one one one two and two two two one, all vanish.}}
\PSlide{Slide 54}{The $2\times2$ CI problem}{Slide 54: The 2×2 CI problem}
\PLine{By symmetry, the \textbf{open-shell singlet} sigma g sigma u is alone in its symmetry, while the two \textbf{closed shells}, sigma g squared and sigma u squared, have the same symmetry and mix.}
\PLine{Their \textbf{Hamiltonian} is a two by two matrix. The diagonal holds the two \textbf{closed-shell} energies, and the first one, E one, is the restricted Hartree Fock energy. The coupling is K one two, the closed-shell to closed-shell element we derived earlier.}
\PLine{The eigenvalues satisfy the determinant condition: E one minus E, times E two minus E, minus K one two squared, equals zero.}
\PLine{Multiplied out, that is a quadratic in E.}
\PLine{Its lower root is the \textbf{configuration interaction} ground state: one half of E one plus E two, minus the square root of the gap squared plus four K one two squared.}
\PLine{The \textbf{triplet} and the open-shell singlet of sigma g sigma u keep their earlier forms: h one one plus h two two plus J one two, minus or plus K one two.}
\PSlide{Slide 55}{Pulling the atoms apart (1/2): the integrals}{Slide 55: Pulling the atoms apart (1/2): the integrals}
\PLine{Now pull the two atoms apart, to infinite separation. The overlap goes to zero, and so does every overlap density, chi A times chi B. Call U the repulsion of two electrons on the same atom, and V the repulsion of electrons on different atoms, which tends to one over R.}
\PLine{The bonding orbital squared contains a cross term, two chi A chi B, which vanishes: the density becomes one half chi A squared plus chi B squared, half an electron on each atom.}
\PLine{The antibonding orbital squared tends to the same thing.}
\PLine{The product of the two orbitals is one half of chi A plus chi B times chi A minus chi B: one half chi A squared minus chi B squared.}
\PLine{Every \textbf{Coulomb integral} is then the repulsion of two half-clouds, spread over both atoms: four cross terms, U plus V plus V plus U, over four. One half of U plus V.}
\PLine{The \textbf{exchange integral} uses the product density, where chi B squared enters with a minus sign: U minus V minus V plus U, over four. One half of U minus V.}
\PLine{And each one-electron integral tends to the hydrogen atom energy minus one over R, because an electron on atom A still feels nucleus B as a point charge.}
\PSlide{Slide 56}{Pulling the atoms apart (2/2): the energies}{Slide 56: Pulling the atoms apart (2/2): the energies}
\PLine{Now the energies, including the nuclear repulsion, one over R. Restricted Hartree Fock gives two h one one plus J one one plus one over R: two hydrogen energies, minus two over R, plus U over two plus V over two, plus one over R. With V equal to one over R, that is two E H plus U over two minus one over two R. It does not reach two separate hydrogen atoms: it is too high by U over two.}
\PLine{The \textbf{triplet} gives two E H minus two over R, plus V, plus one over R, which is exactly two E H: the correct limit.}
\PLine{In \textbf{configuration interaction}, E one and E two become equal, so the lower root is simply E one minus K one two. Subtracting U minus V over two cancels the error: the result tends to two E H.}
\PLine{The \textbf{open-shell singlet} lies two K above the triplet: two E H plus U minus V, an \textbf{ionic} state.}
\PLine{Why does Hartree Fock fail? Half of sigma g squared describes H plus and H minus, an ion pair, which costs U over two. In configuration interaction the two \textbf{closed shells} are degenerate, and the coupling K one two cancels exactly that cost.}
\PLine{\Take{So restricted Hartree Fock misses the dissociation limit by U over two, and one extra configuration fixes it exactly.}}
\PSlide{Slide 57}{What the CI wavefunction is}{Slide 57: What the CI wavefunction is}
\PLine{At infinite separation, the \textbf{configuration interaction} eigenvector has c one equal to minus c two, one over root two each: two configurations with equal weight.}
\PLine{Expand sigma g squared in atomic orbitals. It contains \textbf{covalent} terms, one electron on each atom, and \textbf{ionic} terms, both electrons on the same atom, with equal weight.}
\PLine{Now take the configuration interaction combination, sigma g squared minus sigma u squared.}
\PLine{The ionic terms cancel, and what remains is purely covalent: chi A of one chi B of two plus chi B of one chi A of two.}
\PLine{\Take{So Hartree Fock is fifty percent ionic, and the ion pair at infinity costs U over two, while configuration interaction is purely covalent: the Heitler London wavefunction, one hydrogen atom on each side.}}
\PSlide{Slide 58}{H$_2$ curves from our calculation}{Slide 58: H2 curves from our calculation}
\PLine{Here are the curves from our own \textbf{PySCF} calculation of H two in the minimal STO three G basis. The dashed line is two separate hydrogen atoms, and the gold curve is restricted Hartree Fock.}
\PLine{The two-configuration result, which here equals \textbf{full configuration interaction}, bends over and approaches the dashed line.}
\PLine{The \textbf{triplet} is repulsive and also reaches two E H, and the \textbf{open-shell singlet} lies above it.}
\PLine{\textbf{Configuration interaction} and the triplet converge on two separate atoms, but Hartree Fock does not: test T five finds it still zero point three six two hartree too high at twenty bohr, exactly K one two. The closed forms derived above agree with PySCF to four times ten to the minus sixteen hartree.}
\PSlide{Slide 59}{The general rule: orbital gap versus exchange}{Slide 59: The general rule: orbital gap versus exchange}
\PLine{The same mechanism works for any two orbitals. Take a model with orbital energies minus and plus delta over two, same-orbital repulsion J d, and different-orbital repulsion J x.}
\PLine{The two \textbf{closed shells} have energies J d minus delta and J d plus delta, coupled by K.}
\PLine{The lower root is J d minus the square root of delta squared plus K squared.}
\PLine{From the first row of the eigenvalue equation, the ratio of the two coefficients is delta minus that square root, over K.}
\PLine{Plot it against delta over K. When the gap is large, the ratio goes to zero: one configuration, a \textbf{closed shell}. When the gap closes, the ratio goes to minus one: two configurations of equal weight, a pure \textbf{diradical}. Twisting ethylene to ninety degrees closes the gap between pi and pi star, which is the case of Figure two b.}
\PLine{The \textbf{singlet} \textbf{triplet} gap becomes J d minus J x, plus K, minus that square root.}
\PLine{\Take{So the weaker the interaction, the stronger the \textbf{static correlation}, exactly as the review says.} Test T six checked the model symbolically.}
\PPart{Part IX: Kinetic exchange}{Part IX: Kinetic exchange}
\PSlide{Slide 60}{The Hubbard picture}{Slide 60: The Hubbard picture}
\PLine{The next mechanism is \textbf{kinetic exchange}. The review says: the kinetic exchange picture based on the \textbf{Hubbard} model attributes the stability of a \textbf{singlet} state, for orthogonal atomic orbitals, to the fact that the perturbation yielding an \textbf{ionic} configuration is allowed only from the singlet, and not from the \textbf{triplet}.}
\PLine{Set it up with two orthonormal orbitals, l and r, localised on the left and right radical centres.}
\PLine{An electron can hop between them, with matrix element minus t.}
\PLine{Putting two electrons on the same site costs the on-site repulsion U, two electrons on different sites repel by V, and there is a direct \textbf{exchange integral k} between the sites, which is never negative.}
\PSlide{Slide 61}{Step 1: ionic configurations are pure singlets}{Slide 61: Step 1: ionic configurations are pure singlets}
\PLine{\textbf{Step one.} In an \textbf{ionic} configuration both electrons sit in one orbital, for example l of one times l of two.}
\PLine{That spatial part is symmetric: swapping the electrons changes nothing.}
\PLine{Since the total must be \textbf{antisymmetric}, its spin part must be the antisymmetric \textbf{singlet}.}
\PLine{The \textbf{Hamiltonian} contains no spin, so any matrix element between an ionic configuration and a \textbf{triplet} is proportional to the overlap of a singlet spin function with a triplet spin function, which is zero. Hopping an electron from one site to the other can connect a \textbf{covalent} singlet to an ionic one, but never a triplet.}
\PLine{\Take{So only the singlet can borrow energy from ionic configurations.}}
\PSlide{Slide 62}{Step 2: the matrix}{Slide 62: Step 2: the matrix}
\PLine{\textbf{Step two.} Insert these parameters into the two-electron two-orbital \textbf{Hamiltonian} from Part five: the site energies are zero, and the mixed integral l l l r vanishes.}
\PLine{The \textbf{triplet} energy is V minus k. The singlets form a three by three matrix in the two \textbf{ionic} configurations and the \textbf{covalent} one: U on the ionic diagonal, V plus k on the covalent diagonal, k between the two ionic configurations, and minus root two t between each ionic configuration and the covalent one.}
\PLine{Every entry is one we derived before, with the new names.}
\PLine{\textbf{Step three.} Combine the two ionic configurations into a symmetric and an \textbf{antisymmetric} pair, Phi plus and Phi minus.}
\PLine{Their energies are U plus k and U minus k.}
\PLine{Their couplings to the covalent configuration add for Phi plus, giving minus two t, and cancel for Phi minus, giving zero.}
\PLine{\Take{So Phi minus decouples, and the ground \textbf{singlet} lives in the two by two space of the covalent configuration and Phi plus.}}
\PSlide{Slide 63}{Step 4: solve and expand}{Slide 63: Step 4: solve and expand}
\PLine{\textbf{Step four.} Solve the two by two problem. Its lower root is one half of U plus V plus two k, minus the square root of U minus V squared plus sixteen t squared.}
\PLine{Subtract the \textbf{triplet} energy, V minus k: the gap is one half U minus V, plus two k, minus one half of that square root.}
\PLine{Now expand the square root for x equal to U minus V, much larger than t: it is x plus eight t squared over x plus smaller terms. The x terms cancel, leaving two k minus four t squared over U minus V.}
\PLine{Two contributions compete. Plus two k is potential exchange, which favours the triplet. Minus four t squared over U minus V is \textbf{kinetic exchange}: the energy the \textbf{singlet} gains by virtually hopping into the \textbf{ionic} configuration, which the triplet cannot do. It favours the singlet.}
\PSlide{Slide 64}{Hubbard dimer: exact versus kinetic exchange}{Slide 64: Hubbard dimer: exact versus kinetic exchange}
\PLine{The simplest case is the \textbf{Hubbard} dimer: set V and k to zero. Then the exact \textbf{singlet} energy is one half of U minus the square root of U squared plus sixteen t squared, and the \textbf{triplet} energy is zero. Here is the exact gap, in units of U, against t over U.}
\PLine{The \textbf{kinetic exchange} approximation, minus four t squared over U, agrees with it when t is much smaller than U, which is the regime of strong correlation.}
\PLine{In that regime the two electrons behave like two spins coupled by a Heisenberg \textbf{Hamiltonian}, J A F times S one dot S two, with J A F equal to four t squared over U. S one dot S two is minus three quarters in the singlet and plus one quarter in the triplet, so the singlet is lower by exactly J A F.}
\PLine{Applied bond by bond, this coupling will build the topology rules later on.}
\PLine{Test T seven checked the whole derivation symbolically and against \textbf{PySCF} \textbf{full configuration interaction}, with t equal to one, U six, V one and a half and k zero point two.}
\PSlide{Slide 65}{Ionic resonance reverses a spin state: Figure 6}{Slide 65: Ionic resonance reverses a spin state: Figure 6}
\PLine{The review’s Figure six shows \textbf{ionic} resonance reversing a spin state. Meta xylylene, a classic \textbf{triplet} \textbf{diradical}, sits next to its pyridinium analog, molecule five.}
\PLine{The review says: in the pyridinium analog of meta xylylene, this contribution, resonance with \textbf{closed-shell} zwitterionic configurations, exceeds the exchange and Coulomb terms, which formally favour the triplet of this system with \textbf{nondisjoint} non-bonding orbitals.}
\PLine{In our model, low-lying zwitterionic structures mean a small effective U minus V. Then the \textbf{kinetic exchange} term, four t squared over U minus V, grows past two k, and the \textbf{singlet} wins.}
\PLine{That mapping is qualitative: the review gives no numbers to test.}
\PSlide{Slide 66}{Delocalised and localised: two views of one state}{Slide 66: Delocalised and localised: two views of one state}
\PLine{Finally, these two pictures are the same physics. Build delocalised orbitals a and b as the sum and difference of l and r. Then a of one a of two minus b of one b of two equals l of one r of two plus r of one l of two.}
\PLine{Expand it: the l l and r r terms, the \textbf{ionic} ones, cancel, and the cross terms add.}
\PLine{In delocalised orbitals, the state looks like two \textbf{closed-shell} configurations of equal weight: \textbf{static correlation}. In localised orbitals, it is a single \textbf{covalent} configuration with no ionic part.}
\PLine{It is the same state with different bookkeeping, so the static-correlation and kinetic-exchange descriptions agree.}
\PLine{One remark from the derivations. The \textbf{Hubbard} model has no separate kinetic and potential energy, so it cannot be held to the \textbf{virial theorem}; like the frozen-orbital model of \textbf{Hund’s rule}, it cannot say which physical energy component pays for the spin preference.}
\PPart{Part X: NBMOs \& disjointness}{Part X: NBMOs \& disjointness}
\PSlide{Slide 67}{Borden \& Davidson’s classification}{Slide 67: Borden \& Davidson’s classification}
\PLine{Part ten turns to the molecules themselves. The review’s Figure three collects high-spin and low-spin polyradicals: cyclobutadiene, tetramethyleneethane, trimethylenemethane, meta xylylene and trimethylenebenzene.}
\PLine{Their spin follows Borden and Davidson’s classification. The review says: if the two orbitals do not have common atoms and the probability of two electrons occupying the same space is nil independent of their spins, they are termed \textbf{disjoint}. Under such circumstances, the energy of an \textbf{open-shell singlet} can become comparable to or even lower than that of a \textbf{triplet}.}
\PLine{The two orbitals are the \textbf{non-bonding molecular orbitals}, or \textbf{N B M Os}, whose energy equals that of an isolated p orbital. Three things need deriving: first, how to find them; second, when they can be disjoint; and third, why disjoint means the \textbf{exchange integral K} is close to zero.}
\PSlide{Slide 68}{Hückel theory and the zero-sum rule}{Slide 68: Hückel theory and the zero-sum rule}
\PLine{To find \textbf{N B M Os} we use \textbf{Hückel} theory. Its assumptions: one p pi orbital per carbon, a diagonal energy alpha, a coupling beta between bonded neighbours and zero otherwise, and zero overlap. \Take{So the \textbf{Hamiltonian} matrix is alpha times the identity plus beta times the adjacency matrix A.} Careful: alpha and beta here are Hückel energies, not spin functions.}
\PLine{Write each orbital energy as alpha plus x beta.}
\PLine{Subtract alpha times c from both sides and divide by beta: the coefficient vector is an eigenvector of the adjacency matrix, with eigenvalue x.}
\PLine{A non-bonding orbital has energy exactly alpha, so x is zero: A times c is the zero vector.}
\PLine{Row i of the adjacency matrix has a one for each neighbour of atom i, so row i of this equation says the coefficients on the neighbours of atom i add up to zero.}
\PLine{That is the \textbf{zero-sum rule}: for every atom, the \textbf{N B M O} coefficients of its neighbours sum to zero.}
\PLine{The allyl radical shows it. The end atoms carry plus and minus one over root two, and the centre carries zero. The centre’s neighbours are the two ends, and their coefficients cancel; each end’s only neighbour is the centre, whose coefficient is zero. Every atom is satisfied.}
\PSlide{Slide 69}{Alternant hydrocarbons: the pairing theorem}{Slide 69: Alternant hydrocarbons: the pairing theorem}
\PLine{Most of these molecules are \textbf{alternant} hydrocarbons: you can star every other atom so that bonds only join starred atoms to unstarred ones. That is possible exactly when there are no odd rings.}
\PLine{Order the starred atoms first. Then the adjacency matrix has zero diagonal blocks and an off-diagonal block B with its transpose.}
\PLine{The zero blocks say that no bond joins two starred atoms or two unstarred atoms.}
\PLine{Split an eigenvector into u on the starred atoms and v on the unstarred ones. Multiplying the blocks, the eigenvalue equation becomes two equations: B v equals x u, and B transpose u equals x v.}
\PLine{Now flip the sign of v. A times u, minus v, is minus B v on top and B transpose u below, which is minus x u and x v, which is minus x times the flipped vector.}
\PLine{\Take{So every level x has a partner at minus x: the levels come in pairs mirrored about alpha.} This is the \textbf{pairing theorem}, and our \textbf{Lean} file proves it as the theorem named pairing.}
\PSlide{Slide 70}{How many NBMOs?}{Slide 70: How many NBMOs?}
\PLine{How many \textbf{N B M Os} does an \textbf{alternant} hydrocarbon have? Each nonzero level of B gives a plus and minus x pair of levels of A, so the rank of A is twice the rank of B. The rank counts the nonzero eigenvalues.}
\PLine{A rank can never exceed the smaller side of the matrix, so it is at most twice the smaller of n A and n B.}
\PLine{The number of N B M Os is the number of zero eigenvalues, N minus the rank. With N equal to n A plus n B, that is at least the difference between the two set sizes.}
\PLine{A sharper bound, due to Longuet-Higgins: the rank of B is at most T, the size of a maximum matching, meaning the largest set of bonds that share no atom. \Take{So the number of N B M Os is at least N minus two T.}}
\PLine{Careful: these are only lower bounds. Square cyclobutadiene has four atoms and a perfect matching of two bonds, so N minus two T is zero, yet it has two N B M Os.}
\PLine{We had first assumed the bound was an equality; test T eight caught that \Err{wrong} assumption.}
\PSlide{Slide 71}{Tetramethyleneethane: solving the zero-sum equations}{Slide 71: Tetramethyleneethane: solving the zero-sum equations}
\PLine{Tetramethyleneethane, T M E, has six pi atoms: a central bond from atom one to atom two, with atoms three and four on atom one, and five and six on atom two. Apply the \textbf{zero-sum rule} atom by atom.}
\PLine{Atom three has a single neighbour, atom one, so its rule forces c one to be zero.}
\PLine{Likewise atom five’s only neighbour is atom two, so c two is zero. Atoms four and six give the same two conditions again.}
\PLine{Atom one’s neighbours are two, three and four. With c two zero, c four equals minus c three.}
\PLine{Atom two’s neighbours are one, five and six. With c one zero, c six equals minus c five.}
\PLine{Two coefficients remain free, c three and c five, so there are two \textbf{N B M Os}. Choosing one free parameter at a time and normalising gives phi a, chi three minus chi four over root two, and phi b, chi five minus chi six over root two.}
\PLine{They share no atom: the N B M Os of T M E are \textbf{disjoint}.}
\PSlide{Slide 72}{Disjoint pairs, drawn by our JavaScript}{Slide 72: Disjoint pairs, drawn by our JavaScript}
\PLine{Our JavaScript draws these orbitals. Lobe area is proportional to the squared coefficient, rose for plus and blue for minus. On the left, T M E: phi a lives on atoms three and four, phi b on five and six.}
\PLine{On the right, square cyclobutadiene: its two \textbf{N B M Os} can also be chosen to share no atom, each living on one diagonal pair.}
\PLine{The pipeline: build the \textbf{Hückel} matrix, diagonalise it by Jacobi rotations, then rotate the degenerate pair to minimise the sum over atoms of a i squared times b i squared, and write S V G, then P D F. A degenerate pair can be mixed freely, so we pick the mixture that overlaps least; when that sum is zero, no atom carries both orbitals.}
\PSlide{Slide 73}{Trimethylenemethane cannot be made disjoint}{Slide 73: Trimethylenemethane cannot be made disjoint}
\PLine{Trimethylenemethane, T M M, is different. Number the centre zero and the terminals one to three. Each terminal’s only neighbour is the centre, so the \textbf{zero-sum rule} there forces c zero to be zero.}
\PLine{At the centre, the rule says c one plus c two plus c three is zero. \Take{So both \textbf{N B M Os} live in this two-dimensional plane of terminal coefficients.}}
\PLine{A vector in this plane with only one nonzero entry cannot sum to zero, so every \textbf{N B M O} has at least two nonzero coefficients.}
\PLine{Two such vectors with no atom in common would need at least four terminal atoms, and there are only three.}
\PLine{\Take{So T M M is \textbf{nondisjoint}, for every choice of basis: no rotation of its two N B M Os can separate them, as the drawing shows.}}
\PLine{Meta xylylene is also nondisjoint.}
\PSlide{Slide 74}{Why disjoint means $K_{ab}=0$}{Slide 74: Why disjoint means Kab=0}
\PLine{Why does disjointness kill the \textbf{exchange integral}? Assume zero differential overlap: any product of two different atomic orbitals is treated as zero, so the two-electron integral over atomic orbitals mu nu lambda sigma survives only when mu equals nu and lambda equals sigma, and then it is gamma mu lambda, the repulsion between an electron on atom mu and one on atom lambda. Expand both orbitals in atomic orbitals and insert them into K a b: a quadruple sum.}
\PLine{Zero differential overlap keeps only the terms with nu equal to mu and sigma equal to lambda: K becomes a double sum of a mu b mu, times a lambda b lambda, times gamma.}
\PLine{If the orbitals are \textbf{disjoint}, the product a mu b mu vanishes on every atom, so every term vanishes and K is zero. Conversely, gamma is positive definite, so for \textbf{nondisjoint} orbitals K is strictly positive.}
\PLine{In the \textbf{Hubbard} model, where only on-site repulsion U survives, K is U times the sum over atoms of a squared b squared.}
\PLine{Minimised over the choice of basis, that sum is one sixth for T M M and one sixteenth for meta xylylene; our Python and JavaScript agree on both.}
\PSlide{Slide 75}{The test: how the gap grows with $U$}{Slide 75: The test: how the gap grows with U}
\PLine{Test T nine b checks all of this with exact diagonalisation of the \textbf{Hubbard} model, every determinant and no approximation, at t equal to one. Both axes are logarithmic: the singlet-triplet gap against U over t. On log-log axes a slope of one means the gap grows like U, which is first order, two K; a slope of two means it grows like U squared, which is second order.}
\PLine{The \textbf{nondisjoint} molecules, T M M and meta xylylene, have \textbf{triplet} ground states, and their gaps rise with slope about one. At U equal to zero point zero five, T M M’s gap is zero point zero one six four, against the first-order prediction two U times one sixth, zero point zero one six seven.}
\PLine{The \textbf{disjoint} molecules, cyclobutadiene and T M E, behave differently: their gaps rise with slope about two, fitted at one point nine seven and one point nine five, and the \textbf{singlet} lies below the triplet.}
\PLine{\Take{So the first-order \textbf{exchange integral} really does vanish for disjoint orbitals, and something at second order picks the singlet.} The next part finds out what.}
\PPart{Part XI: Spin polarisation}{Part XI: Spin polarisation}
\PSlide{Slide 76}{Dynamic spin polarisation}{Slide 76: Dynamic spin polarisation}
\PLine{Part eleven: \textbf{spin polarisation}, and the topology rules it produces. The review says: a damped spin wave so formed generates ferromagnetic coupling between the spins separated by an odd number of sites, and antiferromagnetic between those separated by an even number of sites, thus yielding the topology rules of spin coupling. Ferromagnetic means parallel spins; antiferromagnetic means antiparallel; and sites are atoms.}
\PLine{Here is the mechanism in three steps. First, in \textbf{Hückel} theory the unpaired electron sits in the non-bonding orbital, which is exactly zero on the unstarred atoms.}
\PLine{Second, switch on the on-site repulsion U, which acts only between an up and a down electron on the same atom. Where the unpaired up-spin electron sits, down-spin density is pushed away, onto the neighbouring atoms.}
\PLine{Third, negative spin density appears on the unstarred atoms, between the positive ones. Hückel theory cannot produce it: it is a correlation effect.}
\PSlide{Slide 77}{Measured on the pentadienyl radical}{Slide 77: Measured on the pentadienyl radical}
\PLine{Test T nine measured this on the pentadienyl radical, a chain of five carbons with one unpaired electron, by exact diagonalisation of the \textbf{Hubbard} model: five sites, five electrons, three up and two down. At U equal to zero, which is \textbf{Hückel} theory, the spin density is one third on the starred sites one, three and five, and exactly zero on sites two and four.}
\PLine{At U equal to t, small negative dips appear on the unstarred sites.}
\PLine{At U equal to three t, they reach minus zero point one eight on sites two and four.}
\PLine{At six t, minus zero point two five: the damped spin wave of the review.}
\PLine{At every U the total stays one, positive on the starred atoms and negative on the unstarred ones.}
\PSlide{Slide 78}{From bonds to topology}{Slide 78: From bonds to topology}
\PLine{How does a coupling on each bond become a rule about topology? First, at strong U each bond carries the antiferromagnetic coupling J A F, four t squared over U: \textbf{kinetic exchange}, from Part nine.}
\PLine{Second, an \textbf{alternant} molecule is a bipartite graph: every bond joins a starred atom to an unstarred one.}
\PLine{Third, the lowest-energy arrangement of the spins puts all starred spins one way and all unstarred spins the other, so every bond is antiparallel. This is the Néel arrangement.}
\PLine{Fourth, two atoms are in the same set exactly when an even number of bonds separates them, which means an odd number of atoms in between.}
\PLine{\Take{So: same set, parallel; different sets, antiparallel.} In meta xylylene the two C H two groups have three ring atoms between them, so their spins are parallel and the ground state is a \textbf{triplet}. In para xylylene there are four atoms between them: antiparallel, a \textbf{singlet}.}
\PSlide{Slide 79}{Ovchinnikov’s rule, and Lieb’s theorem behind it}{Slide 79: Ovchinnikov’s rule, and Lieb’s theorem behind it}
\PLine{Count the Néel arrangement: n A spins up on the starred atoms and n B down on the unstarred ones. The total spin is one half of the difference, n A minus n B.}
\PLine{That is \textbf{Ovchinnikov}’s rule. As the review states it, the most spin-alternant determinant is expected to have the lowest energy and the highest contribution.}
\PLine{Behind it stands a theorem, by \textbf{Lieb}, in Physical Review Letters, nineteen eighty nine: for the repulsive \textbf{Hubbard} model on a bipartite graph at half filling, with any U greater than zero, the ground state is unique and has spin one half of the difference between the sizes of the two sets.}
\PLine{Test T nine checked the rule against exact diagonalisation at U equal to three t, for seven molecules: cyclobutadiene, T M M, T M E, meta xylylene, trimethylenebenzene, tetramethylenebenzene and para xylylene. In every case the exact spin equals the rule’s count.}
\PLine{One correction to the review. On page four hundred and forty four it groups trimethylenebenzene with meta xylylene as having \textbf{triplet} ground states. Its three unpaired electrons give a quartet, with spin three halves, as both the rule and our calculation say.}
\PSlide{Slide 80}{Disjoint diradicals: who picks the singlet?}{Slide 80: Disjoint diradicals: who picks the singlet?}
\PLine{Back to the question left open in Part ten: for \textbf{disjoint} \textbf{diradicals}, what picks the \textbf{singlet}? The review’s answer, for square cyclobutadiene: in the singlet state, the bonding pi electrons can partially localise on the same pair of atoms as the same-spin unpaired electrons; in the \textbf{triplet} state such advantageous \textbf{spin polarisation} is not possible.}
\PLine{At first order the gap is two K a b, which is zero for disjoint orbitals.}
\PLine{At second order the paired electrons polarise, mixing in spin-polarised excited configurations, and that energy scales as U squared.}
\PLine{Our scaling test found exactly that: the gap grows as U to the power one point nine seven for cyclobutadiene and one point nine five for T M E, with the singlet lower.}
\PLine{That is consistent with the spin-polarisation mechanism of the review’s reference nine, by Kollmar and Staemmler, titled Violation of \textbf{Hund’s rule} by spin polarization, which we could not access.}
\PSlide{Slide 81}{Topology in the paper: Figure 4}{Slide 81: Topology in the paper: Figure 4}
\PLine{The review’s Figure four shows topology at work in real molecules.}
\PLine{It says that simple polyradicals can also be thought of as coupling units that sustain certain spin alignment depending on topology, as seen in Schlenk’s and Chichibabin’s hydrocarbons, or in isomers of dinitroxide four.}
\PLine{In the figure, dinitroxide four with para, para prime links has parallel spins, while the meta, meta prime and meta, para prime isomers have antiparallel spins. The atoms are the same; only the topology differs.}
\PSlide{Slide 82}{Tropylidene: a singlet from spin polarisation, Figure 5}{Slide 82: Tropylidene: a singlet from spin polarisation, Figure 5}
\PLine{\textbf{Spin polarisation} can even produce a \textbf{singlet} ground state where \textbf{Hund’s rule} would predict a \textbf{triplet}. The review’s Figure five shows tropylidene, a carbene in a seven-membered ring. Follow the mechanism in five steps.}
\PLine{One: an unpaired electron occupies the two a two pi orbital.}
\PLine{Two: it polarises the bonding pi electrons in two b one.}
\PLine{Three: as a result, opposite spin accumulates on the carbenic carbon.}
\PLine{Four: that opposite spin couples ferromagnetically, following Hund’s rule, with the electron in the orthogonal in-plane a one orbital on the same atom.}
\PLine{Five: the net result is that the singlet, one A two, lies lower than the triplet.}
\PLine{Our summary: Hund’s rule is obeyed locally, on the carbenic carbon, and violated globally, for the molecule as a whole.}
\PSlide{Slide 83}{Where the rules stop}{Slide 83: Where the rules stop}
\PLine{Where do these rules stop working? \textbf{Ovchinnikov}’s rule and \textbf{Lieb}’s theorem cover \textbf{alternant} pi systems, where \textbf{Hubbard} or Heisenberg type physics applies.}
\PLine{The pyridinium molecule five, from Figure six, is an exception: \textbf{ionic} resonance overturns the topology prediction.}
\PLine{And sigma pi \textbf{diradicals}, such as the didehydrotoluenes, break the pattern too: the ortho and para isomers are triplets, but the meta isomer is an \textbf{open-shell singlet}.}
\PLine{The review itself says: crucially, all of these rules are not the direct consequences of quantum mechanics and only mimic the dynamic and \textbf{static correlation} effects with the help of simplifying assumptions, thus inevitably having limited scope.}
\PLine{That is fair. The topology rule is a theorem only of the Hubbard model, by Lieb; the others are model heuristics.}
\PPart{Part XII: Switching spin states}{Part XII: Switching spin states}
\PSlide{Slide 84}{Spin-state crossover by external stimuli: Figure 7}{Slide 84: Spin-state crossover by external stimuli: Figure 7}
\PLine{Part twelve: switching spin states. The review’s Figure seven collects molecules whose spin state can be switched by an external stimulus.}
\PLine{In panel a, oxidation: a quinone coupler switches two nitroxide radicals between ferromagnetic and antiferromagnetic coupling.}
\PLine{In panel b, light: trans azobenzene couples its spins antiparallel, a \textbf{singlet}; light turns it into the cis isomer, which loses planarity and is predicted, computationally, to be a \textbf{triplet}.}
\PLine{In panel c, light again: a trans to cis switch of an alkene brings two TEMPO units close together, hydrogen bonded, and they couple antiferromagnetically. TEMPO is a stable nitroxide radical.}
\PLine{In panel d, light forms a new bond, and that changes the orbital topology.}
\PLine{In panel e, zinc ions: coordination creates oxygen zinc oxygen \textbf{superexchange} and through-space exchange between the oxygens, which favour antiferromagnetic coupling. Superexchange is coupling of two spins through an atom in between.}
\PLine{And in panel f, temperature or a magnetic field turns one four phenylenedinitrene from an \textbf{open-shell singlet} into a triplet.}
\PSlide{Slide 85}{A wording slip in (e)?}{Slide 85: A wording slip in (e)?}
\PLine{Panel e comes with a sentence worth reading closely. The review says: coordination to zinc brings about oxygen zinc oxygen \textbf{superexchange} and through-space exchange between the oxygens, which favour antiferromagnetic coupling and override the preference for the \textbf{triplet} ground state, predicted on the basis of \textbf{disjoint} \textbf{N B M Os} in the parent dinitroxide ligand.}
\PLine{But on page four hundred and forty three, the review itself says that disjoint N B M Os make the \textbf{singlet} competitive, and its triplet examples, meta xylylene and trimethylenebenzene, are \textbf{nondisjoint}.}
\PLine{\Take{So disjoint here probably means nondisjoint, or else the triplet preference has some other origin.}}
\PLine{We mark it inconsistent with page four hundred and forty three. We did not check the original, reference fifty two.}
\PSlide{Slide 86}{Panel (f), quantitatively: temperature}{Slide 86: Panel (f), quantitatively: temperature}
\PLine{Panel f can be made quantitative. Assume a \textbf{singlet} ground state with the \textbf{triplet} an energy delta above it, and Boltzmann statistics. The singlet has degeneracy one and the triplet three, one for each M S. \Take{So the partition function is one plus three e to the minus delta over k B T.}}
\PLine{The fraction of molecules in the triplet is its weight over the partition function: three e to the minus delta over k B T, divided by one plus three e to the minus delta over k B T.}
\PLine{At high temperature the exponential tends to one, so the fraction tends to three over one plus three: three quarters.}
\PLine{That agrees with the review: increased temperature promotes electrons to populate the triplet state.}
\PSlide{Slide 87}{Panel (f), quantitatively: magnetic field}{Slide 87: Panel (f), quantitatively: magnetic field}
\PLine{Now a magnetic field. Assume a Zeeman term, g mu B times B times M S, with g about two. The \textbf{singlet} has only M S equal to zero, so it does not move. The three \textbf{triplet} components split apart: delta plus g mu B B times M S.}
\PLine{The M S equal to minus one component reaches the singlet when delta minus g mu B B is zero, at the field B star equal to delta over g mu B.}
\PLine{That matches the review: an external magnetic field splits the triplet into microstates, one of which approaches the energy of the ground-state singlet.}
\PSlide{Slide 88}{Closed shell or diradical? Figures 8 and 9}{Slide 88: Closed shell or diradical? Figures 8 and 9}
\PLine{Some molecules sit on the border between a \textbf{closed shell} and a \textbf{diradical}. In Figure eight, one three dinitrophenyls with anionic substituents transfer an electron within the molecule to give anionic \textbf{diradicals}.}
\PLine{Their \textbf{triplet} is the ground state, or lies only negligibly above the closed shell; and binding a cation to the nitro group stabilises the diradical.}
\PLine{In Figure nine, \textbf{protonation} and oxidation create diradicals: by donor acceptor electron transfer after double protonation, in a sapphyrin-derived \textbf{singlet} biradicaloid, in a dipyrylium diradical salt, and in high-spin building blocks for polaronic ferromagnets.}
\PLine{The review also mentions singlet fission: one excited singlet shares its excitation with a neighbour, and both end in a low-lying triplet.}
\PSlide{Slide 89}{Protonation-induced magnetisation: Figure 9}{Slide 89: Protonation-induced magnetisation: Figure 9}
\PLine{Here is Figure nine in full: protonation-induced magnetisation. The arrows are reactions, adding or removing protons and transferring electrons, and the dots are unpaired electrons, coloured by spin.}
\PLine{This part of the review is descriptive. There is no mathematics to derive, only claims about specific molecules, which we did not re-check.}
\PPart{Part XIII: Koopmans \& Janak}{Part XIII: Koopmans \& Janak}
\PSlide{Slide 90}{The claim}{Slide 90: The claim}
\PLine{Part thirteen: \textbf{Koopmans} and \textbf{Janak}. The review mentions an alternative interpretation of \textbf{Hund’s rule}: an alternative theoretical interpretation of the \textbf{Hund’s multiplicity rule} has been proposed on the basis of Janak theorem, or extended Koopmans’ theorem: at the Hartree Fock level, the singlet-triplet gap is equal to the difference between the energies of the \textbf{triplet}’s and \textbf{open-shell singlet}’s highest occupied orbitals.}
\PLine{We derive it in three steps: Koopmans’ theorem, then Janak’s theorem, then the gap.}
\PLine{Which difference, and with which sign? The review does not say. The derivation will.}
\PSlide{Slide 91}{Koopmans’ theorem}{Slide 91: Koopmans’ theorem}
\PLine{\textbf{Koopmans}’ theorem first. Assume a single determinant of spin orbitals chi i, and remove the electron in chi k without relaxing any of the others.}
\PLine{By the Slater Condon rules, the energy of the N electron determinant is the sum of the one-electron energies, plus one half the double sum over pairs of Coulomb minus exchange integrals.}
\PLine{The determinant without chi k has the same sums with k left out. \Take{So the difference keeps only the terms that contain k: its one-electron energy, the half sum with i equal to k, and the half sum with j equal to k.} The term with both equal to k is counted twice, but it is k k k k minus k k k k, which is zero.}
\PLine{The two half sums are equal, because the integrals are symmetric under swapping their electron pairs, so together they make one full sum of Coulomb minus exchange between k and every j.}
\PLine{That is exactly the definition of the Fock eigenvalue epsilon k.}
\PLine{\Take{So the ionisation potential of orbital k is minus epsilon k, with \textbf{frozen orbitals}.}}
\PSlide{Slide 92}{Janak’s theorem}{Slide 92: Janak’s theorem}
\PLine{\textbf{Janak}’s theorem works in Kohn Sham \textbf{density functional} theory with fractional occupations n i. The energy is the occupation-weighted kinetic energy, plus the nuclear attraction of the density, plus the Hartree and exchange correlation energies, where the density is the sum of n i times the orbital density.}
\PLine{Differentiate with respect to n i. The explicit dependence: n i enters the density linearly, so by the chain rule the derivative is the expectation value of the Kohn Sham operator in orbital i, which is epsilon i.}
\PLine{The orbitals also change with n i. But the energy’s derivative with respect to each orbital is n j epsilon j times that orbital, so this contribution is proportional to the change in each orbital’s norm, which stays one. It vanishes.}
\PLine{\Take{So the derivative of the energy with respect to an occupation is the orbital energy, for any exchange correlation functional.} This is Janak’s theorem, Physical Review B, nineteen seventy eight.}
\PLine{With integer occupations and \textbf{frozen orbitals}, this slope becomes a finite difference, and \textbf{Koopmans}’ theorem comes back.}
\PSlide{Slide 93}{The gap as a HOMO difference}{Slide 93: The gap as a HOMO difference}
\PLine{Now the gap. Define the \textbf{HOMO} energy of a state as minus the cost of removing its phi b electron, with phi a frozen. What remains is the ion, one electron in phi a, with energy h a. For the \textbf{triplet}, the HOMO energy is the triplet energy minus the ion’s: h b plus J minus K, using the triplet energy from Part five.}
\PLine{For the \textbf{open-shell singlet}, the same subtraction gives h b plus J plus K.}
\PLine{Subtract: the ion’s energy cancels, and the singlet-triplet gap, E O S S minus E T, equals the \textbf{singlet}’s HOMO energy minus the triplet’s, which is two K. \Take{So the review’s statement holds, with this sign: the singlet’s HOMO lies higher.}}
\PLine{With \textbf{frozen orbitals}, the triplet’s HOMO is lower, or harder, by two K, in Pearson’s sense. The review adds the caveat, from its reference twenty, that this lacks dynamic electron correlation; orbital relaxation is missing too. Verified, within those limits.}
\PPart{Part XIV: Aufbau on trial}{Part XIV: Aufbau on trial}
\PSlide{Slide 94}{The prosecution}{Slide 94: The prosecution}
\PLine{Part fourteen puts the \textbf{aufbau} principle on trial. Aufbau is German for building up: fill orbitals in order of energy, two electrons each. The prosecution, in the review’s words: the aufbau principle, by assuming the existence of unchanging stationary individual electronic states, violates the \textbf{Pauli} exclusion principle. Only the many-electron system as a whole possesses stationary states.}
\PLine{We examine three exhibits. \textbf{Exhibit A}: one-electron quantum numbers are not conserved.}
\PLine{\textbf{Exhibit B}: the adiabatic argument behind building up.}
\PLine{\textbf{Exhibit C}: the paradox of the three d and four s orbitals.}
\PSlide{Slide 95}{Exhibit A: $[\hat l_{1z},1/r_{12}]\neq0$}{Slide 95: Exhibit A: [l1z,1/r12]≠0}
\PLine{\textbf{Exhibit A}. The z angular momentum of electron one alone is minus i times x one d by d y one minus y one d by d x one. By the chain rule on one over r one two, its derivatives with respect to y one and x one are minus y one minus y two over r cubed, and minus x one minus x two over r cubed.}
\PLine{Apply l one z to one over r one two by inserting both derivatives.}
\PLine{Multiply out the numerator: minus x one y one, plus x one y two, plus y one x one, minus y one x two.}
\PLine{The first and third terms cancel.}
\PLine{Only the cross terms survive: minus i times x one y two minus x two y one, over r cubed. That is not zero.}
\PLine{Because l one z is a first-order differential operator, its commutator with the repulsion acting on any function is just l one z applied to the repulsion, times the function. It is not zero, so the angular momentum of one electron is not a constant of motion.}
\PSlide{Slide 96}{… but the total is conserved}{Slide 96: … but the total is conserved}
\PLine{But the total is conserved. For electron two the same steps give the opposite sign, because the derivative with respect to y two of one over r one two has a plus sign.}
\PLine{Add the two: the cross terms are equal and opposite, and the sum is zero.}
\PLine{\Take{So in an atom, with a central field and a spin-free \textbf{Hamiltonian}, the total orbital angular momentum and the total spin squared commute with the Hamiltonian, but the angular momentum of each electron does not.}}
\PLine{The reason is geometric. The total angular momentum rotates both electrons together and keeps their distance; the angular momentum of electron one rotates it alone and changes the distance. In a molecule even the total orbital angular momentum is lost, while the total spin survives.}
\PLine{Test T thirteen checked these commutators in \textbf{SymPy} on an arbitrary function.}
\PSlide{Slide 97}{Exhibit B: the adiabatic theorem}{Slide 97: Exhibit B: the adiabatic theorem}
\PLine{\textbf{Exhibit B}. The review says that Bohr assumed stationary states survive the addition of an electron and a proton; that this rests on Ehrenfest’s adiabatic principle; and that the principle does not apply to aperiodic systems.}
\PLine{What does the adiabatic theorem actually say? Expand the state in the eigenstates of the \textbf{Hamiltonian} at each instant, with amplitudes c n and dynamic phases theta n, minus the integral of the energy.}
\PLine{Insert this into the time-dependent Schrödinger equation and project on psi m. The phase terms cancel the energies, and what is left is: the rate of change of c m is minus the sum over n of c n times the overlap of psi m with the time derivative of psi n, times a phase factor.}
\PLine{Differentiate the eigenvalue equation and project on psi m: that overlap is the matrix element of the time derivative of the Hamiltonian, divided by the energy gap.}
\PLine{\Take{So if the Hamiltonian changes slowly and the gaps stay finite, the couplings are small, and the system stays in the continuously connected eigenstate.} But adding a particle changes the Hilbert space itself, so there is no continuous path; and the labels that are conserved are L and S of the whole system, never one-electron labels.}
\PSlide{Slide 98}{Exhibit C: the 3d/4s paradox}{Slide 98: Exhibit C: the 3d/4s paradox}
\PLine{\textbf{Exhibit C}, the three d four s paradox. The review says: in neutral atoms four s lies slightly above three d when both are populated, so four s is ionised first; however, four s is also populated first, because it offers less repulsion with a three d electron than that between two three d electrons; orbital energies must be re-evaluated in the new configurations. Let us make this precise.}
\PLine{Use the average-of-configuration model. Each orbital type has a one-electron energy, h d or h s, and each type of pair has one average repulsion, J bar d d, J bar d s, and J s s. Each electron pays its h once, and each pair pays its J once; n electrons form n times n minus one over two pairs.}
\PLine{\Take{So a configuration with n d electrons in three d and n s in four s has energy n d h d plus n s h s, plus n d choose two times J bar d d, plus n d times n s times J bar d s, plus n s choose two times J s s.}}
\PSlide{Slide 99}{3d/4s: the algebra}{Slide 99: 3d/4s: the algebra}
\PLine{Now the algebra. With n three d electrons and two four s electrons: n h d plus two h s, plus n choose two J bar d d, plus two n J bar d s, plus one J s s.}
\PLine{Move one four s electron into three d: n plus one h d plus h s, plus n plus one choose two J bar d d, plus n plus one J bar d s, and no s s pair.}
\PLine{Subtract, calling the difference delta. Since n plus one choose two minus n choose two is n, delta is h d minus h s, plus n J bar d d, plus one minus n times J bar d s, minus J s s.}
\PLine{The orbital energies in d n s two, as frozen-orbital removal energies: removing a three d electron costs h d plus n minus one J bar d d plus two J bar d s, and removing a four s electron costs h s plus n J bar d s plus J s s.}
\PLine{Subtract the orbital energies from delta: every h and every J s s cancels, the J bar d d terms leave one, the J bar d s terms leave minus one. \Take{So delta equals epsilon d minus epsilon s, plus J bar d d minus J bar d s.} Test T twelve checks this symbolically, and the \textbf{Lean} theorem three d four s proves it.}
\PSlide{Slide 100}{3d/4s: reading the identity}{Slide 100: 3d/4s: reading the identity}
\PLine{How to read the identity. A compact three d orbital and a diffuse four s orbital suggest that J bar d d is larger than J bar d s; that is a heuristic, not something we derived.}
\PLine{Then even when epsilon d lies below epsilon s, delta is positive if the repulsion difference exceeds the orbital energy gap. The configuration d n s two wins: four s fills first.}
\PLine{But starting from d n s two, \textbf{Koopmans}’ theorem makes the four s ionisation energy the smaller one, so four s leaves first.}
\PLine{Some illustrative numbers, with n equal to two: epsilon s minus zero point nine, epsilon d minus one point five, J bar d d zero point eight and J bar d s zero point one hartree. The orbital gap gives minus zero point six, the repulsion difference gives plus zero point seven, and delta is plus zero point one: d two s two is lower, although three d is the lower orbital.}
\PLine{The moral: orbital energies belong to a configuration, so the order of building up and the order of removal need not agree. There is no paradox.}
\PSlide{Slide 101}{Verdict on the aufbau critique}{Slide 101: Verdict on the aufbau critique}
\PLine{\Take{The verdict.} What the review gets right: one-electron quantum numbers are labels of a model, as \textbf{Exhibit A} showed;}
\PLine{orbital energies depend on the configuration, as \textbf{Exhibit C} showed;}
\PLine{and \textbf{non-aufbau} occupations are observed through measurable properties: in the vanadyl porphyrin V O o e p, in quantum dots, in the dissociation of lithium fluoride, and in the radicals coming up next.}
\PLine{What is overstated: a determinant filled by \textbf{aufbau} is fully \textbf{antisymmetric}, so within the orbital model it does not violate \textbf{Pauli}’s principle.}
\PLine{And the obsolescence of aufbau and \textbf{Hund’s rule} is too strong: both remain useful rules of thumb for atomic ground states, with known exceptions such as chromium and copper.}
\PLine{Overall verdict on the aufbau critique: overstated.}
\PPart{Part XV: Quasi-closed shells}{Part XV: Quasi-closed shells}
\PSlide{Slide 102}{SOMO below HOMO: Figure 10}{Slide 102: SOMO below HOMO: Figure 10}
\PLine{Part fifteen: \textbf{quasi-closed} shells, or \textbf{non-aufbau} radicals. In an ordinary radical the singly occupied orbital, the \textbf{SOMO}, is the highest occupied orbital. The review’s Figure ten shows the alternative: the SOMO lies below one or more doubly occupied orbitals.}
\PLine{The review says: upon ionisation an electron is removed from the higher energy doubly occupied orbital, leaving another unpaired electron behind, while upon addition the new electron enters the former \textbf{LUMO}; essentially, the SOMO does not participate in these processes.}
\PLine{\Take{So the fingerprint is \textbf{redox} chemistry: oxidation and reduction give polyradicals, \textbf{diradicals} in the simplest case, not \textbf{closed shells}.}}
\PSlide{Slide 103}{UHF orbital energies of $s^1h^2$, part 1}{Slide 103: UHF orbital energies of s1h², part 1}
\PLine{Why should a \textbf{SOMO} ever drop below a \textbf{HOMO}? Derive the orbital energies. Assume \textbf{frozen orbitals}: a SOMO s with one alpha electron, a doubly occupied orbital h, the core absorbed into the one-electron energies, and each orbital energy defined, as in \textbf{Koopmans}’ theorem, as the cost of removing that electron.}
\PLine{The configuration is s alpha, h alpha, h beta. There are three pairs: s alpha with h alpha, which have the same spin; s alpha with h beta; and h alpha with h beta.}
\PLine{The energy: h s plus two h h, plus, for each pair, a \textbf{Coulomb integral}, and for the same-spin pair also minus the \textbf{exchange integral K} h s.}
\PLine{Remove h alpha: it takes its one-electron energy and its two pairs with it. \Take{So epsilon h alpha is h h plus J h h plus J h s minus K h s.}}
\PLine{Remove h beta instead: it has no same-spin partner, so there is no K. Epsilon h beta is h h plus J h h plus J h s.}
\PSlide{Slide 104}{UHF orbital energies of $s^1h^2$, part 2}{Slide 104: UHF orbital energies of s1h², part 2}
\PLine{Now the \textbf{SOMO} itself. Removing s alpha breaks its pair with h alpha, worth J minus K, and its pair with h beta, worth J. \Take{So epsilon s alpha is h s plus two J h s minus K h s.}}
\PLine{The beta SOMO is empty, so its energy is the cost of adding an electron there. The new electron pairs with s alpha, paying the large on-site repulsion J s s, with h alpha, paying J h s, and with h beta, which has the same spin, paying J h s minus K h s. \Take{So epsilon s beta is h s plus J s s plus two J h s minus K h s.}}
\PLine{Compare the two \textbf{HOMO} energies: epsilon h beta minus epsilon h alpha is K h s, which is never negative. The donor’s alpha HOMO is pulled down by exchange with the unpaired electron; its beta HOMO is not. This is the review’s Scheme three, derived.}
\PLine{Test T two checked all four energies against energies of explicit determinants.}
\PSlide{Slide 105}{When does the SOMO drop below the HOMO?}{Slide 105: When does the SOMO drop below the HOMO?}
\PLine{When does the \textbf{SOMO} drop below the \textbf{HOMO}? Compare the two alpha energies. Epsilon s alpha is below epsilon h alpha exactly when h s plus two J h s minus K h s is less than h h plus J h h plus J h s minus K h s.}
\PLine{The exchange terms cancel, and so does one J h s: the condition is h s plus J h s less than h h plus J h h.}
\PLine{Read it: on the left, the radical orbital’s energy plus its repulsion with one h electron; on the right, h’s own energy plus its self-repulsion.}
\PLine{\Take{So conversion is favoured by a stable radical with a low h s, next to an electron-rich fragment, for example an anionic one, with a high h h.}}
\PLine{Exactly the review’s recipe: a spatially separated stable radical centre and a molecular fragment serving as a source of high-energy doubly occupied orbitals.}
\PSlide{Slide 106}{The redox fingerprint, derived}{Slide 106: The redox fingerprint, derived}
\PLine{Now derive the \textbf{redox} fingerprint. Oxidation first: by \textbf{Koopmans}’ theorem the electron leaves from the orbital with the smallest ionisation potential, which is h, since minus epsilon h is less than the \textbf{SOMO}’s.}
\PLine{The product, s one h one, is a \textbf{diradical}.}
\PLine{Its states are those of Part five: the \textbf{open-shell singlet} lies two K s h above the \textbf{triplet}.}
\PLine{\Take{So oxidation gives triplet cation \textbf{diradicals} with nondegenerate SOMOs, as observed.}}
\PLine{Reduction: the new electron goes either into s beta, at epsilon s beta, or into the former \textbf{LUMO} l, at epsilon l.}
\PLine{At \textbf{frozen orbitals}, the product is the diradical s one h two l one if and only if epsilon l is below h s plus J s s plus two J h s minus K h s.}
\PLine{That is a separate condition: SOMO \textbf{HOMO} conversion on its own does not imply it.}
\PLine{What makes it possible is the large on-site repulsion J s s, which pushes epsilon s beta up.}
\PSlide{Slide 107}{ROHF: are these orbital energies even real?}{Slide 107: ROHF: are these orbital energies even real?}
\PLine{Before trusting these orbital energies, a warning from the review: erroneous \textbf{non-aufbau} occupations can be observed with \textbf{restricted open-shell} methods, depending on the canonicalisation parameters of the alpha and beta electrons in the Fock matrix. Let us see why.}
\PLine{In restricted open-shell Hartree Fock there are three spaces: the doubly occupied core c, the \textbf{open shell} o, occupied by alpha electrons only, and the virtuals v. The alpha density covers core and open shell; the beta density covers the core.}
\PLine{A small change of the densities changes the energy by the trace of each Fock matrix times the change of its density. Rotating orbitals p and q into each other by a small angle kappa changes a density by kappa times the symmetric combination of p and q.}
\PLine{But a rotation changes a spin density only if exactly one of p and q is occupied in that spin. Two occupied orbitals, or two empty ones, can be mixed without changing anything.}
\PSlide{Slide 108}{ROHF stationarity, rotation by rotation}{Slide 108: ROHF stationarity, rotation by rotation}
\PLine{Go through the rotations one by one, asking for each whether it changes the alpha density, the beta density, and so what stationarity requires.}
\PLine{Core with \textbf{open shell}: both are alpha-occupied, so the alpha density does not change, but the beta density does. Stationarity requires the beta Fock element between core and open shell to vanish.}
\PLine{Open shell with virtual: the alpha density changes, the beta density does not, since both are beta-empty. \Take{So the alpha Fock element between open and virtual must vanish.}}
\PLine{Core with virtual: both densities change, so the sum of the alpha and beta elements must vanish.}
\PLine{And rotations within the core, within the open shell, or within the virtuals change nothing at all: the energy is blind to them, and there is no condition.}
\PLine{\Take{So the diagonal blocks of the Fock matrix are arbitrary.} Any weighted combination, A times the alpha Fock matrix plus one minus A times the beta one, gives the same energy and the same density, but different orbital energies.}
\PLine{Test T eleven a confirmed it for H two N O, N O two, methyl and allyl: the required blocks are below ten to the minus seven, while the others are about zero point one. \textbf{PySCF}’s Roothaan Fock matrix uses exactly these off-diagonal blocks.}
\PSlide{Slide 109}{How arbitrary? H₂NO, ROHF/6-31G∗}{Slide 109: How arbitrary? H₂NO, ROHF/6-31G∗}
\PLine{How arbitrary is it in practice? We took H two N O, the smallest nitroxide, with \textbf{restricted open-shell} Hartree Fock in the six thirty one G star basis, and varied the weight A in the open-open block of the effective Fock matrix. The energy is identical for every A: minus one hundred and thirty point three six nine seven zero nine hartree.}
\PLine{The \textbf{HOMO}, a \textbf{closed-shell} orbital, stays put at minus zero point four six six seven. But the \textbf{SOMO} energy moves linearly with A, from plus zero point two at A equal to zero to minus zero point four four at A equal to one.}
\PLine{The SOMO crosses the HOMO at A star, about one point zero four.}
\PLine{The natural range of weights, from zero to one, is shaded: there is no crossing inside it.}
\PLine{\Take{So the review is right in principle, but for ordinary radicals the standard choices do not flip the order.} We mark it partial.}
\PSlide{Slide 110}{Hence: trust observables}{Slide 110: Hence: trust observables}
\PLine{\Take{Hence the review’s advice, which we endorse: a rigorous support for the \textbf{non-aufbau} occupation can be obtained after the fact from the ionisation potentials and electron affinities, or \textbf{redox} potentials in solution, as well as from the spin state and magnetic properties of the open-shell oxidation and reduction products.}}
\PLine{\Take{In other words, stop arguing about orbital ordering, which depends on a convention, and compare the total energies of the \textbf{diradical} and the \textbf{closed shell}, which do not.}}
\PLine{That is exactly the test we ran on a \textbf{distonic} radical anion, a few slides ahead.}
\PSlide{Slide 111}{Inorganic precedents: Scheme 2}{Slide 111: Inorganic precedents: Scheme 2}
\PLine{\textbf{Non-aufbau} radicals were known first in inorganic chemistry. The review’s Scheme two shows the vanadyl porphyrin V O o e p. Its unpaired electron is localised on vanadium.}
\PLine{Yet one-electron oxidation removes an electron from an occupied porphyrin orbital,}
\PLine{and one-electron reduction puts it into an unoccupied porphyrin orbital. The vanadium electron takes part in neither.}
\PLine{The same behaviour has been reported in quantum dots and along the dissociation curve of lithium fluoride.}
\PSlide{Slide 112}{A TEMPO-dithiolate on platinum: Figure 11}{Slide 112: A TEMPO-dithiolate on platinum: Figure 11}
\PLine{An organometallic example, the review’s Figure eleven: a TEMPO dithiolate ligand on platinum. Oxidising the complex removes an electron from the dithiolate \textbf{HOMO}, not from the TEMPO \textbf{SOMO}.}
\PLine{The evidence: U V visible spectra, electron spin resonance, cyclic voltammetry, \textbf{density functional} calculations, and the product of dimerisation.}
\PLine{The free ligand is a normal radical. The conversion comes from an antibonding interaction of its next-highest occupied orbital with the HOMO of the platinum bipyridine fragment, which pushes it above the SOMO.}
\PLine{In panel b, a ditempo metalladithiolene: oxidation generates biradicals that couple to give a bicyclic product.}
\PSlide{Slide 113}{Spin-polarised donors: Figures 12 and 13}{Slide 113: Spin-polarised donors: Figures 12 and 13}
\PLine{Figure twelve: nitronyl nitroxide radicals attached to donors such as tetrathiafulvalene, benzoquinone and amines. One-electron oxidation gives \textbf{triplet} cation \textbf{diradicals}, the fingerprint we derived.}
\PLine{Figure thirteen links the effect to disjointness, from Part ten.}
\PLine{With a one-pyrrolyl link, the donor \textbf{HOMO} is \textbf{disjoint} from the \textbf{SOMO} and the next orbital. With a three-pyrrolyl link, the donor HOMO and the next-highest orbital are \textbf{nondisjoint}, and their combination rises above the SOMO.}
\PSlide{Slide 114}{Conducting magnets: Scheme 4, Figure 14, Scheme 5}{Slide 114: Conducting magnets: Scheme 4, Figure 14, Scheme 5}
\PLine{The same orbital pattern gives conducting magnets, Scheme four. After oxidation there are two nondegenerate SOMOs: the donor’s, higher, carries an itinerant electron that can conduct, and the radical’s, lower, carries a localised spin that gives magnetism.}
\PLine{Scheme five: the perchlorate salts of E T B N and E S B N, described as the first genuine paramagnetic conductors.}
\PLine{Their neutral crystals show negative magnetoresistance and have been used in field-effect transistors.}
\PLine{And Figure fourteen: several radicals on one donor; oxidation gives a strongly ferromagnetic, high-spin system.}
\PSlide{Slide 115}{Other radicals, and a pH sensor: Figure 15}{Slide 115: Other radicals, and a pH sensor: Figure 15}
\PLine{Figure fifteen extends the idea to other radicals. In panel a, a nitronyl nitroxide carries a phenol-based donor. Removing or adding a proton switches the oxidised molecule between a \textbf{triplet} and a \textbf{closed-shell} \textbf{singlet}: a pH sensor.}
\PLine{In panel b, an anthracene, para phenylene, oxoverdazyl molecule gives high-spin oxidation products.}
\PSlide{Slide 116}{Distonic radical anions: Figure 16}{Slide 116: Distonic radical anions: Figure 16}
\PLine{Figure sixteen shows \textbf{distonic} radical anions: the charge and the radical sit on different atoms.}
\PLine{The review’s claim: a stabilised radical plus an anionic group not conjugated to it; the anion’s \textbf{HOMO} lies above the \textbf{SOMO}; if the anion is protonated, the \textbf{aufbau} occupation is restored, so it can be switched by changing the pH; and on oxidation it has been shown computationally to give \textbf{diradicals} in preference to the \textbf{closed-shell} zwitterion.}
\PLine{That is testable, and so far only computational. \Take{So we computed both oxidised states, with and without the proton.}}
\PSlide{Slide 117}{Our test: two radicals, two protonation states, two functionals}{Slide 117: Our test: two radicals, two protonation states, two functionals}
\PLine{Our test, T eleven and T eleven c. Two radicals, a short one and a TEMPO-based one, each as the acid and as the anion; geometries from \textbf{unrestricted} B three L Y P; and vertical oxidation computed with two functionals. The bars show the energy of the oxidised \textbf{triplet} \textbf{diradical} minus that of the oxidised \textbf{closed shell}. Negative means oxidation gives the diradical; positive, the closed shell.}
\PLine{With B three L Y P, the acids oxidise to \textbf{closed shells}, by one point four and one point three electron volts, and in them the \textbf{SOMO} is the \textbf{HOMO}.}
\PLine{With the range-separated functional CAM B three L Y P, the anions oxidise to \textbf{diradicals}, by zero point three six and zero point three zero electron volts: the pH switch, reproduced.}
\PLine{But with B three L Y P the anions miss, by six to seven hundredths of an electron volt on the \textbf{closed-shell} side.}
\PSlide{Slide 118}{What the test says}{Slide 118: What the test says}
\PLine{What did the test show? Reproduced: in the anions the \textbf{SOMO} lies two to three levels below the \textbf{HOMO}, and in the acids it is the HOMO.}
\PLine{\textbf{Protonation} moves the \textbf{triplet} minus \textbf{closed-shell} energy by about one point two to one point eight electron volts.}
\PLine{And with CAM B three L Y P the anions oxidise to \textbf{diradicals}, as the review says.}
\PLine{The caveat: for the short anion, B three L Y P lowers the \textbf{closed shell} by zero point six one electron volts relative to CAM B three L Y P, but the triplet by only zero point one seven.}
\PLine{\Take{The verdict flips at about seven hundredths of an electron volt: the claim is sensitive to the functional.}}
\PLine{Delocalisation error, the tendency of approximate functionals to over-stabilise charge spread over many atoms, is a plausible cause; we did not test it.}
\PSlide{Slide 119}{Non-aufbau by electron addition: Figure 17}{Slide 119: Non-aufbau by electron addition: Figure 17}
\PLine{Finally, \textbf{non-aufbau} behaviour on electron addition, Figure seventeen. In panel a, the spirobiphenalenyl radical six is reduced to an \textbf{open-shell singlet}, and it crystallises without dimerising: an intrinsic molecular metal.}
\PLine{In panel b, galvinoxyl added to a P three H T and P C B M solar cell: its \textbf{LUMO} mediates the crossing of polaron pairs from \textbf{singlet} to \textbf{triplet} through exchange, which prevents recombination and raises the photovoltaic efficiency.}
\PLine{And self-doping: in a tetrathiafulvalene carboxylate ammonium salt, the \textbf{SOMO} of the oxidised tetrathiafulvalene unit is not the highest occupied orbital, according to calculations on a model tetramer.}
\PPart{Part XVI: Broken symmetry}{Part XVI: Broken symmetry}
\PSlide{Slide 120}{Cheap methods for diradicals}{Slide 120: Cheap methods for diradicals}
\PLine{Part sixteen: \textbf{broken symmetry}. The review says that relatively cheap methods, such as \textbf{unrestricted} or \textbf{broken-symmetry} formalisms, suffer from a number of serious drawbacks, while rigorous methods are too costly for large systems.}
\PLine{A broken-symmetry determinant puts the alpha electron in an orbital a prime on the left centre and the beta electron in b prime on the right, with overlap s.}
\PLine{For any unrestricted determinant, the expectation value of S squared is S z times S z plus one, plus the number of beta electrons, minus the sum of squared overlaps between alpha and beta orbitals.}
\PLine{For the broken-symmetry determinant, S z is zero, there is one beta electron, and one overlap s: the result is one minus s squared. This formula is exact for a determinant.}
\PLine{If s is one, the two orbitals coincide: a \textbf{closed shell}, with S squared zero. If s is zero, the determinant is half \textbf{singlet} and half \textbf{triplet}, with S squared one.}
\PLine{As Part four showed, a single determinant cannot be the \textbf{open-shell singlet}.}
\PSlide{Slide 121}{Yamaguchi’s projection, derived}{Slide 121: Yamaguchi’s projection, derived}
\PLine{How do \textbf{broken-symmetry} calculations still give couplings? Through \textbf{Yamaguchi}’s projection, which we derive. Assume the two-centre Heisenberg model, minus two J times S one dot S two, and write the broken-symmetry state as a mixture of the \textbf{singlet} and the M S zero \textbf{triplet}.}
\PLine{\textbf{Step one.} Since S squared is S one squared plus S two squared plus twice S one dot S two, and each S i squared is three quarters, S one dot S two is one half of S times S plus one, minus three halves. \Take{So the triplet energy is minus J over two and the singlet energy three J over two.}}
\PLine{\textbf{Step two.} S squared is zero in the singlet and two in the triplet, and the two states do not mix under either operator, so the broken-symmetry state has S squared equal to two c T squared, and energy c S squared E S plus c T squared E T.}
\PLine{\textbf{Step three.} Subtract the triplet energy: c S squared times E S minus E T. With c T squared equal to one minus c S squared, and E S minus E T equal to two J, this is one minus half the broken-symmetry S squared, times two J.}
\PLine{\textbf{Step four.} Solve for J: the energy difference between the broken-symmetry and high-spin determinants, divided by the difference of their S squared values, two for the high-spin triplet.}
\PLine{Our \textbf{Lean} file proves this as the theorem yamaguchi. Sign conventions differ between sources; here a negative J means antiferromagnetic coupling.}
\PSlide{Slide 122}{How good is it? The Hubbard dimer}{Slide 122: How good is it? The Hubbard dimer}
\PLine{How good is the projection? Test T fourteen compares it with the exact answer for the \textbf{Hubbard} dimer. The exact coupling, one half of E S minus E T, is one quarter of U minus the square root of U squared plus sixteen t squared.}
\PLine{The dashed curve is \textbf{Yamaguchi}’s formula applied to self-consistent \textbf{broken-symmetry} \textbf{unrestricted} Hartree Fock. At U equal to three t, it underestimates the coupling by seven point seven percent.}
\PLine{At U equal to twelve t, the error is only a tenth of a percent.}
\PLine{\Take{So the projection is good when the centres are weakly coupled, with the broken-symmetry S squared near one, and poor when the bond is \textbf{covalent}.}}
\PLine{The review’s serious drawbacks is a fair description.}
\PPart{Part XVII: Our verification}{Part XVII: Our verification}
\PSlide{Slide 123}{The repository}{Slide 123: The repository}
\PLine{Part seventeen: our verification. Everything is in one repository. The paper itself; eighteen LaTeX sections of derivations; the verification scripts, tests one to fifteen, with their logs and results; the \textbf{Lean} proofs; the review’s figures and captions; the slide deck; documentation of the findings and the test matrix; our working plans and audit; the web page; and this video, with its scenes, narration and editor.}
\PLine{Each artefact is one command. make sync builds the locked environment with uv; make verify runs tests one to fifteen into the results folder; make verify d f t runs the \textbf{density functional} tests; make lean checks the proofs; and make figures regenerates the figures byte for byte.}
\PLine{Python three point one two is pinned by the lock file, with \textbf{PySCF} two point one four and \textbf{SymPy} one point one four; and Lean four point two nine point one with the matching Mathlib.}
\PSlide{Slide 124}{The test matrix, part 1}{Slide 124: The test matrix, part 1}
\PLine{Here is the test matrix. Each test names what it checks, how, and the result.}
\PLine{T one checks the spin functions and the combinations of Figure one, including the sign \Err{erratum}, in \textbf{SymPy}. T one b checks the sign \Err{erratum} again, numerically, by applying S squared to wavefunction tensors in NumPy.}
\PLine{T two compares the two-electron two-orbital \textbf{Hamiltonian} with \textbf{PySCF} \textbf{full configuration interaction} for five random integral sets. T three checks that no \textbf{exchange integral} is negative and none exceeds the \textbf{Coulomb integral}, for water.}
\PLine{T four solves helium’s one s two s states exactly and checks the \textbf{virial theorem} and \textbf{less screening}. T five checks the dissociation limits of H two.}
\PLine{T six checks the two-orbital gap model, and T seven \textbf{kinetic exchange} in the extended \textbf{Hubbard} dimer, symbolically and against full configuration interaction.}
\PLine{T eight checks the \textbf{Hückel} rules for non-bonding orbitals and disjointness, and T nine the spins predicted by \textbf{Ovchinnikov} and \textbf{Lieb}, and spin alternation, by exact diagonalisation. All pass.}
\PSlide{Slide 125}{The test matrix, part 2}{Slide 125: The test matrix, part 2}
\PLine{The second half. T nine b measures how the gap grows with U. T ten checks that the singlet-triplet gap is a \textbf{HOMO} difference.}
\PLine{T eleven a checks the stationarity conditions of \textbf{restricted open-shell} Hartree Fock. T eleven b, the \textbf{distonic} anion with B three L Y P, fails, and we report it as it is. T eleven c repeats it for two radicals, two \textbf{protonation} states and two functionals, with its own checking script, and passes.}
\PLine{T twelve checks the three d four s identity, T thirteen the angular momentum commutators, T fourteen the \textbf{Yamaguchi} projection, and T fifteen same-spin against opposite-spin correlation.}
\PLine{T seventeen is the \textbf{Lean} file: seventeen algebraic theorems, machine-checked.}
\PLine{T eleven b fails as recorded: B three L Y P misses the crossover by six to seven hundredths of an electron volt, and T eleven c shows that the outcome depends on the functional. T sixteen, the citations, were checked by hand against publisher records, so it has no script.}
\PSlide{Slide 126}{Inside T4: helium, exact in a 16s7p4d basis}{Slide 126: Inside T4: helium, exact in a 16s7p4d basis}
\PLine{Inside test T four. This is the actual code, unedited. The two-electron wavefunction is expanded in pairs of basis functions, phi p of electron one times phi q of electron two, with a coefficient matrix C.}
\PLine{The \textbf{Hamiltonian} in this pair space is h for electron one, h for electron two, and the repulsion integral, assembled with einsum and reshaped into a matrix.}
\PLine{For the \textbf{singlet}, C must be symmetric; for the \textbf{triplet}, \textbf{antisymmetric}, because the spin factor carries the other symmetry. The loop builds a basis of symmetric or antisymmetric pairs.}
\PLine{The Hamiltonian projected onto that basis is diagonalised. For the singlet the second root is the one s two s state; for the triplet it is the lowest.}
\PLine{Then the kinetic, nuclear attraction and repulsion energies are computed from C, and the code asserts that they add up to the eigenvalue, to ten to the minus eight. These are the numbers behind the virial analysis of \textbf{Hund’s rule}.}
\PSlide{Slide 127}{Inside T11a: ROHF Fock blocks}{Slide 127: Inside T11a: ROHF Fock blocks}
\PLine{Inside test T eleven a. Run \textbf{restricted open-shell} Hartree Fock to tight convergence, and take its density matrices.}
\PLine{Build the alpha and beta Fock matrices from those densities, as \textbf{unrestricted} Hartree Fock would, and transform them to the molecular orbital basis. Then sort the orbitals into core, open and virtual by their occupations.}
\PLine{The blocks the stationarity table says must vanish, beta core open, alpha open virtual, and the alpha plus beta core virtual sum, must be below ten to the minus five; while the blocks the table leaves free must be clearly non-zero, above ten to the minus three, to show the test can tell them apart.}
\PSlide{Slide 128}{Inside T17: Lean proves the less-screening logic}{Slide 128: Inside T17: Lean proves the less-screening logic}
\PLine{Inside test T seventeen, the \textbf{Lean} proof of the less-screening logic. The comment states it: if the \textbf{triplet} is lower, both states obey the \textbf{virial theorem}, and the triplet has the larger repulsion, then the kinetic energy rises, and the nuclear attraction falls by more than twice the energy gap and by more than the repulsion rises.}
\PLine{The theorem takes the four differences as real numbers, with four hypotheses: the virial theorem for the differences, the energy as the sum of its parts, the triplet lower, and the repulsion higher. Its conclusion has four parts, and linarith proves each by linear arithmetic.}
\PLine{The second theorem: \textbf{frozen orbitals}, with no change in kinetic energy, are incompatible with the virial theorem unless the energy gap is zero.}
\PLine{And the axiom audit: for all seventeen theorems, print axioms lists only the three standard axioms of Lean’s logic, propositional extensionality, classical choice and quotient soundness. No sorry, and no extra assumptions.}
\PSlide{Slide 129}{What the tests caught in our own work}{Slide 129: What the tests caught in our own work}
\PLine{The tests also caught our own mistakes. T eight: we had assumed the number of non-bonding orbitals was exactly N minus two T; cyclobutadiene showed it is only a lower bound.}
\PLine{T nine b: the lowest level with S z zero was the M S zero component of the \textbf{triplet}, not a \textbf{singlet}.}
\PLine{T eleven a: we expected the \textbf{SOMO} and \textbf{HOMO} to flip under standard canonicalisation schemes; it needs a weight above one.}
\PLine{T five: a symmetry tolerance of ten to the minus ten failed on round-off, at overlaps near ten to the minus sixteen.}
\PLine{T twelve: our first illustrative numbers had the \Err{wrong} sign of delta.}
\PLine{Every failure was diagnosed before the test was changed, and each is logged in the plans folder.}
\PLine{A check that is able to fail is what turns I believe into checked.}
\PPart{Part XVIII: Verdict}{Part XVIII: Verdict}
\PSlide{Slide 130}{Scorecard, part 1: confirmed}{Slide 130: Scorecard, part 1: confirmed}
\PLine{Part eighteen: the verdict. First, what we confirmed. The \textbf{triplet} is single-reference and the \textbf{open-shell singlet} multi-reference.}
\PLine{\textbf{Less screening}, not a lower electron repulsion, explains \textbf{Hund’s rule}.}
\PLine{\textbf{Static correlation} grows as the interaction between the orbitals weakens.}
\PLine{\textbf{Kinetic exchange}, the \textbf{disjoint} non-bonding orbitals, and the topology rules of \textbf{Ovchinnikov}.}
\PLine{The three d four s argument, and the non-conservation of one-electron quantum numbers.}
\PLine{And \textbf{restricted open-shell} orbital energies do depend on a convention, although the practical danger is overstated.}
